%\chapter{Cutting Planes for LRSP}

\section{Introduction}
In this chapter, we investigate cutting planes valid for LRSP and
that can be incorporated into our branch and price algorithm. Extending
a branch and price algorithm to a branch, price and cut algorithm
using cutting planes is a promising approach in the effort of
improving the overall performance of the solution algorithm. Addition of
cutting planes can improve the LP relaxation bounds in the expense
of time to search for cuts, to solve larger LP models, and possible
performance reduction in the column generation procedure caused by
the incorporation of the cuts. Therefore, efficiency of an extended
algorithm depends on other factors related with branch and price
algorithm besides the strength of the new cutting planes.
%careful (deeper) consideration of

The point of this chapter is to derive some valid inequalities,
include them to our branch and price algorithm and investigate the
performance of the overall algorithm. In order to find cuts that are
valid for LRSP, we can explore the set of valid inequalities derived
for known polytopes that are relaxations and special cases of LRSP,
such as knapsack problem, set partitioning problem, LRP and VRP.


The main contributions of this chapter are to investigate some of
the most common cuts that are valid for the related polytopes of
LRSP; discuss their validity for LRSP, their integration
into the model SPP-LRSP and necessary updates in our branch and
price algorithm; and asses the value of these cuts for the solution
algorithm. In addition, we explore possible ways of adapting the approach of
generating disjunctive cutting planes which are one set of generic
cutting planes that are commonly used in branch and cut algorithms.
There does not exist a current publication that is regarding disjunctive
cuts for branch and price algorithms. This chapter also provides some
theoretical results about the strength of the SPP-LRSP by showing that
some of the cuts that can be derived using edge based formulation are
already implied by the SPP-LRSP.

%Need an intro to this chapter , i.e., why are you generating cuts?
%Conjecture:  (1) Generating cuts to add to the branch-and-price
%algorithm will improve the performance of the algorithm (2) Can
%extend LRP and VRP cuts to LRSP
%Point of chapter 3 is to show that (1) you can extend cuts to LRSP
%and (2) they do or do not improve performance
%What is your contribution ?

%In this chapter we investigate valid inequalities for the LRSP
%polytope.
Next, we review some terminology used in the rest of the chapter and
describe necessary updates to our branch and price algorithm to
extend it to a branch, price and cut algorithm. Then, the analysis in
this chapter can be divided into three parts. First, we
study some relaxations and projections of formulation SPP-LRSP
in order to obtain valid inequalities. Second, we investigate
possible ways to include disjunctive cuts (one of the generic
MIP cuts) in a branch and price algorithm. Third, we consider
inequalities valid for the relaxations of the edge based formulation
of LRSP and discuss their effect to SPP-LRSP formulation.
%that can be incorporated into branch and price algorithm.

\section{Overview of Cut Generation in Branch And Price Algorithm}
In this section, we outline the main steps of a branch, price and
cut algorithm in general. We use the same terminology as in Section
\ref{sec:desribeBP} and the description of the branch and price
algorithm, BRANCH-PRICE, also given in Section \ref{sec:desribeBP}.

First, we summarize some basic terminology from integer programming.
An inequality is called valid inequality if it is satisfied by all points
in the convex hull of a problem. A cut is a valid inequality that is
generally added to the model later in the solution process and possibly
violated by the current LP relaxation solution. One needs to solve the
separation problem to find a cut that is violated by a given fractional
solution. The aim of adding valid inequalities and cuts to the model is
to get a better approximation of the convex hull of the problem.
A valid inequality is called a face if it is satisfied with equality by
some solutions of the problem. A face is called facet if it is nonempty and
its dimension is one less than the dimension of the problem. The most beneficial
valid inequalities are the facets of the convex hull. Therefore,
people take great interest in finding facet defining inequalities.
%has great interest polyhedral study

In this chapter, we search for valid inequalities and cuts for LRSP
and we want to integrate these to our solution algorithm. For ease of
the notation, we refer the procedure that generates cuts for a given
RMP as CUTGEN(RMP) and we refer it as CUTGEN() if the procedure does
not depend on the current LP solution. Both CUTGEN(RMP) and CUTGEN()
return a set $C$ including cuts that are generated at the end of
the procedures. In order to explain the modifications to the
BRANCH-PRICE to include cuts, we classify the generated cutting planes
based on their structure.
\begin{itemize}
\item {\em Global Cut:} If the derived inequality is valid for the
original problem's convex hull, then it is {\em global} to the tree,
i.e.,it is valid for all nodes of the tree.
\item {\em Local Cut:} Instead, if it is valid for the
convex hull of the problem at the specific node, then the cut is
{\em local} to that node, and only valid for that node and its
children.
\item {\em Z-free Type Cut:} If an inequality has zero coefficients
for all variables that are generated via column generation algorithm,
then we refer this cut as {\em z-free type} cut.
\item {\em With-z type Cut:} If some variables from the column generation
have nonzero coefficients in the cut, then we refer it as {\em with-z type}
cut.
\end{itemize}
The incorporation of a cut to the algorithm depends on the structure of
the cut. If a cut is global, then we can add the cut to the RMP in all
active nodes of the branch and price tree. Otherwise, the cut is added to
the RMP at the current node (and at its children). If we need the LP
relaxation solution to generate a violated cut, then the cut is generated
after the column generation algorithm (COLGEN) is terminated and the optimal
RMP solution is found at the current node. Otherwise, the cut may be added at
any time to the current RMP (before or after the COLGEN
or any iteration of the COLGEN). In our algorithm, if the cut
generation function does not need the LP solution, then we add the
valid inequality before we start COLGEN.

Addition of a z-free type cut does not affect the solution algorithm for
the pricing problem since it does not affect the column generation algorithm.
%A z-free type cut can be added to the algorithm at any time during the
%solution of the RMP at the node (before or after the column generation
%or any iteration of the column generation) unless we need the LP
%relaxation solution to generate a violated z-free type cut. If we need
%the LP relaxation solution, we add the cut after the column generation
%algorithm is terminated and the optimal RMP solution is found at the
%current node.
On the other hand, a with-z type cut affects the column
generation algorithm, since it has nonzero coefficients for some
variables that are previously generated or may be generated in the future
via column generation algorithm. Therefore, we need to pass the information
about the cut to the pricing problem. This may require changes in the solution
algorithm for the pricing problem.

We extend the BRANCH-PRICE algorithm to a branch, price and cut algorithm
by updating the solution process at each node of the tree as in the
following:

\fbox{\begin{minipage}{5in}
{\bf A node in BRANCH-PRICE-CUT}
\begin{itemize}
\item[0.] Let RMP be the RMP at the beginning of the node.
For the root node, it is the initial RMP and for a non root node,
it is the parent's final RMP updated with the related branching rule.
\item [1.] Call CUTGEN(). Let $C$ $\leftarrow$ CUTGEN() where $C$ is the
set of cuts.
\item [2.] If $C \neq \emptyset$, then update RMP with cuts, i.e.,
RMP $\leftarrow$ (RMP + $C$).
\item[3.] RMP $\leftarrow$ COLGEN(RMP).
\item[4.] If the RMP solution is infeasible or integer, STOP. Otherwise,
run CUTGEN(RMP). Let $C$ $\leftarrow$ CUTGEN(RMP).
\item[5.] If $C=\emptyset$, STOP processing the node. Otherwise,
update RMP with cuts, RMP $\leftarrow$ (RMP + $C$).
\item[6.]If $C$ has any with-z type cuts, then update COLGEN
based on the information from with-z type cuts.
\item[7.] RMP $\leftarrow$ COLGEN(RMP).
\end{itemize}
\end{minipage}}

%This section should be a description of the branch cut and price solution
%methodology along the lines of section 2.2.  Don't get into
%implementation details; instead keep the discussion at the algorithm
%design level.

\section{Cutting Planes Based on SPP-LRSP}
In this section, we discuss some valid inequalities derived based on
the SPP-LRSP formulation presented in Section \ref{part1:sp_form}.
In Chapter 1, we augmented this formulation with two sets of valid
inequalities, (\ref{lb-open-fac}) and (\ref{relation}), which are
shown to be effective computationally in Section \ref{alternativeForm}.
Here, we study the LRSP polytope, defined as follows, to derive more
valid inequalities. We used the same notation as in Chapter 1.
\begin{optprog}
  & \polyS = \mbox{conv}\{t \in \{0,1\}^{|J|}, z \in \{0,1\}^{\sum_{j \in J}|P_j|}  |
                 t, z \mbox{ satisfies (\ref{spp_con1}), (\ref{spp_con2})}\}.
\end{optprog}
This study is two fold: first, we investigate the cutting planes for
the relaxations of polytope $\polyS$ and second, we investigate the cutting
planes for the projections of $\polyS$.
%and the cutting planes for the projections together with some disjunctive cuts.
%Second, we use fractional solutions for the LP relaxation of SPP-LRS to drive disjunctive cuts.

\subsection{Relaxations of the SPP-LRSP}
 \label{cuts-relax-spp}
 In this section, we investigate two relaxations:
 knapsack and set partitioning problem. We list most common
 set of cuts for these problems (from the literature) based on
 our notation for SPP-LRSP. Even though, we did not implement any of
 these cuts for our branch and price algorithm, we felt the necessity of
 reviewing the literature and discussing possible implementation
 difficulties in the context of this chapter.

\paragraph{Knapsack Polytope}
In $\polyS$, if we relax constraints (\ref{spp_con2}) by dropping
the variable $t$, then each constraint defines a single 0/1 knapsack
polytope. Since we assume that the column set for each facility, in
other words items available for each knapsack in the context of
knapsack problems, are disjoint from each other, we consider $|J|$
single knapsack polytopes instead of one multiple knapsack polytope
in which the set of items available for the knapsacks are not
disjoint.
%then this set of constraints
%become knapsack constraints corresponds to each facility.
We let $\polyKn$ be the knapsack polytope that corresponds to
facility $j \in J$, such that
\begin{optprog}
  & \polyKn = \mbox{conv}\{z \in \{0,1\}^{|P_j|}  | \sum_{p \in P_j} D(p) z_{p} \leq  L^F_{j}\} \ \ & \forall j \in J,
\end{optprog}
%$\mathcal{P}^j_{KNAP1}=\{z_p \in \{0,1\}^{|P_j|} | \sum_{p \in P_j} D^C(p) z_{p} \leq  L^F_{j}\}$
%$\forall j \in J$
where $D(p)$ is equal to the total demand of customers in pairing
$p$ (i.e.,$D(p)= \sum_{i \in I} a_{ip} D_i$).
 $\polyKn$ for all  $j \in J$ is a relaxation of $\polyS$.
%\begin{prop}
%$\mathcal{P}^j_{KNAP1}$ $\forall j \in J$ is a relaxation of
%$\mathcal{P}_{SPP-LRS}$.
%\end{prop}
Thus, we can investigate the valid inequalities of $\polyKn$ which
are also valid for SPP-LRSP.
%improve the LP relaxation of SPP-LRSP.
\citet{Balas75}, \citet{Padberg80}, \citet{Crowder} and
\citet{Weismantel}  are some of the papers studying the polyhedral
structure of knapsack problems. Here, we evaluate applicability of
the following cuts that are investigated in those studies to our algorithm.

One set of well-known nontrivial inequalities that are valid for
$\polyKn$ is the {\em cover inequalities}. Based on our notation and
 definition of $\polyKn$, a set $S_j \subseteq P_j$ is called a cover if
$\sum_{p \in S_j}D(p) > L^F_j$. And if the set $S_j$ is minimal,
then the following inequality is valid for $\polyKn$.
\begin{equation}
\sum_{p \in S_j}z_{p} \leq |S_j|-1 \label{coverinq}
\end{equation}

%The extension of the inequalities (\ref{coverinq}) and the lifting procedures
%to get stronger inequalities are described in \citet{Nemhauser} and \citet{Gu}.
%Lifting is a procedure to extend the inequality (\ref{coverinq}) to also include
%variables $z_p$ for all $p \in P_j \backslash S_j$.

Another set of nontrivial inequalities defined for knapsack
polytopes is {\em (1-k)-configuration inequalities}. Let $S_j$ be a
subset of pairings for which the total demand is no more than the
capacity (i.e.,$S_j \subseteq P_j$ and $\sum_{p \in S_j}D(p)\leq
L^F_j$) and with the property that every one of its subsets is a
minimal cover if you add one pairing from the pairings not in $S_j$.
%$S_j \subseteq P_j$ with $\sum_{p \in S_j}D(p) \leq L^F_j$, and $Q_j
%\cup \{q\}$ is a minimal cover for every $Q_j \subseteq S_j$ where
%$2 \leq |Q_j|=k \leq |S_j|$ and $q \in P_j \backslash S_j$.
Then, the following inequality called (1-k)-configuration inequality
is valid for $\polyKn$.
\begin{equation}
(r - k + 1)z_q + \sum_{p \in S^r_j}z_{p} \leq r, \label{1kconfig}
\end{equation}
where  $r$ is some integer satisfying $k \leq r \leq |S_j|$, and
$S^r_j$ is a subset of $S_j$ such that $|S^r_j| = r$. When $k =
|S_j|$, the inequality (\ref{1kconfig}) is a cover inequality;
therefore cover inequalities are also (1-k)-configuration
inequalities.
%This is a lot of notation to understand!  How about
%adding some interpretation -- In other words, S_j is a subset of
%paths for which the total demand is no more than the capacity and
%with the property that every one of its subsets is a minimal cover
%if you add one element/path from the paths not in S_j.

In a branch and price algorithm, we may separate
%inequalities in terms of
(\ref{coverinq}) and  (\ref{1kconfig}) based on the columns
currently in the model, and we may add these into the current RMP.
However, since these inequalities are with-z type inequalities (with
nonzero coefficients for columns that are generated via column
generation algorithm), unless we incorporate the information from
these generated inequalities into the pricing problem, the
inequalities are not helpful; the pricing problem will generate the
same or similar columns that are visiting the same set of customers.
%However, unless we
%incorporate the information in added inequalities in forms of
%(\ref{coverinq}) and (\ref{1kconfig}) to the pricing problem, the
%pricing problem would generate same or similar columns visiting the
%same set of customers included in previously added inequalities.

Our current solution algorithms for the pricing problem can handle
the new information about a single customer or pair of customers
with strict conditions of either following each other or not
connected to each other in a pairing. In addition, a new
information considering a set of customers can be included in the
pricing algorithm if the information for each customer in the set is
independent of another. These inequalities are constructed
considering pairings and all customers included. The order of each
customer in a pairing depends on the other customers that are
included. Therefore, current pricing algorithms are not equipped to
handle the new information from these generated cuts. At this time,
it is not clear how to update pricing algorithms efficiently
(without making it more difficult to solve) in order to handle
the inequalities in the form of (\ref{coverinq}) and (\ref{1kconfig}).
Therefore we could use neither (\ref{coverinq}) nor (\ref{1kconfig})
in our branch and price algorithm.

Next we discuss the set partitioning polytope and at the end of this
discussion, there is an example from literature, in which the pricing
algorithm is updated based on the new set of cuts under certain conditions.
If there exists an efficient way of handling inequalities in the form
of (\ref{coverinq}) and (\ref{1kconfig}), they would improve the
performance of the algorithm as they are successful in the solution
of $\polyKn$.

%include specific columns and they do not have a structure passing a
%general information about the customers or facilities or the set of
%all columns, thus it is not clear how to define inequalities,
%(\ref{coverinq}) and (\ref{1kconfig}) in pricing problem, and an
%implementation procedure to add these in a branch and price is not
%available at the time. Therefore we could not use them in our branch
%and price algorithm.

%%%%%%%%%%%%%%%%%%%%%%%%%%%%%%%%%%%%%%%%%%%%%%%%%%%%%%%%%%%%%%%%
%%%%%%%%%%%%%%%%%%%%%%%%%%%%%%%%%%%%%%%%%%%%%%%%%%%%%%%%%%%%%%%%%%%
%I think that you can convey this same information in a more positive
%way. What I mean is that, if it is not trivial why one cannot add
%these inequalities, then there is something to be said. See if you
%can better highlight/describe where the difficulty is. It is worth
%mentioning what you anticipate the inequalities would do for you if
%you COULD implement them.  As it stands now, it is not entirely
%clear to the reader WHY you are including the description of these
%inequalities.



\paragraph{Set Partitioning Polytope} Constraints (\ref{spp_con1})
and the integrality constraints in $\polyS$  define set partitioning
polytope. Thus, set partitioning polytope is a relaxation of
$\polyS$, so are set packing and set covering polytopes union of
which results the set partitioning polytope. Cuts for set partitioning
problem are derived by investigating the polyhedral structures of
set packing and set covering polytopes. See \citet{Balas76} for a
review of polyhedral structures of all these polytopes. A common
class of valid inequalities used in the literature is {\em clique
inequalities}. A clique is a complete subgraph. For a set packing or
set partitioning instance one can draw an intersection graph in
which there is one node for each variable and there are edges
between variables if they have positive coefficients in the same
constraint. Based on such graph, the following inequality is called
clique inequality
\begin{equation}
\sum_{p \in S}z_{p} \leq 1 \label{clique-inq}
\end{equation}
where $S$ is the set of nodes in a clique in the intersection graph. Since
all pairings in set $S$ have a coefficient 1 in the same row, we can
only select one of these pairings in an integer solution.
%Lifting procedure is used to strengthen the clique inequalities.

Another common class of inequalities used for set partitioning problems is
{\em odd hole inequalities} that are constructed from cordless cycles with odd
number of vertices in the intersection graphs. Let $C$ be the set of
columns representing vertices in an odd cycle, then the following is valid
for the set partitioning polytope:
\begin{equation}
\sum_{p \in C}z_{p} \leq \Big\lfloor \frac{|C|}{2} \Big\rfloor \label{oddcyc-inq}
\end{equation}

There are other types of inequalities such as anti-hole, webs and
wheels (see \citet{Eso} for details) that are derived from the
intersection graphs in a similar way and have  structure similar
to (\ref{clique-inq}) and (\ref{oddcyc-inq}).
Here, our aim is to provide a short review of cuts valid for set
partitioning problem. Because of the same reason explained for
knapsack cuts earlier in this section, we did not implement any of
these cuts. We do not have a way of updating pricing algorithm
efficiently to incorporate the information from these cuts.
%Here, we do not mention all classes of inequalities and get into
%details of set partitioning, packing and covering polytopes, since
%as in inequalities (\ref{coverinq}) and  (\ref{1kconfig}), we could
%not find a way to incorporate clique (or similar types of)
%inequalities into the pricing problem.

There is a branch, price and cut algorithm example in the literature
(\citet{Jepsen}) in which the authors provide a procedure to update
the pricing problem in order to use the information from a new class
of cuts generated based on the set partitioning polytope. Now, we
review this paper in order to further clarify why updating the
pricing problem for inequalities in the form of (\ref{coverinq}) to
(\ref{oddcyc-inq}) is not trivial and why the structure of the cut is
very important for a branch and price algorithm.
% The problem with inequalities in the form of (\ref{coverinq}) to
%(\ref{oddcyc-inq}) is that they are related with specific columns
%and constructed based on the information about these set of columns.
%To overcome this problem, \citet{Jepsen} proposed {\em subset-row
%inequalities} that are based on the idea of clique and odd hole
%inequalities except that subset-row inequalities are constructed by
%keeping the information from the rows of the set partitioning
%polytope. Based on our notation, the inequalities are
\citet{Jepsen} introduced {\em subset-row inequalities} that are
based on the idea of clique and odd hole inequalities. However, they
constructed these inequalities based on the information from the
rows of the set partitioning polytope instead of considering certain
set of variables as in clique and odd hole inequalities. Based on
our notation, \citet{Jepsen}'s subset-row inequalities are as follows
\begin{equation}
\sum_{j \in J}\sum_{p \in P_j}\Big\lfloor \frac{1}{k} \sum_{i \in I_k}a_{ip} \Big\rfloor z_{p} \leq \Big\lfloor \frac{|I_k|}{k} \Big\rfloor, \label{ssr-inq}
\end{equation}
where $I_k$ is a subset of customers (rows in set partitioning
constraints) and $0 < k \leq |I_k|$. A column visiting at least $k$
of the customers in set $I_k$ has a nonzero coefficient in
inequality (\ref{ssr-inq}), otherwise the coefficient is zero. They
used this condition for a column to update the reduced cost
calculation in pricing problem which is also an ESPPRC.

The reduced cost of a column should be updated with $-\gamma \Big\lfloor \frac{\sum_{i \in
I_k}a_{ip}}{k}\Big\rfloor $ amount where $\gamma$ is the dual variable
that corresponds to constraint (\ref{ssr-inq}).
%They note that this amount is always positive, thus updating
%the cost of a path at the end of the algorithm would be proper. However,
%they use this positive addition as an advantage in
To do that, they kept value $m_l = (|I_k \cap V_l| \mbox{ mod } k)$
for each label $l$ in the ESPPRC algorithm where $V_l$ is the set
of customers in label $l$. They modified domination rules based on
$m_l$ value for the labels. They proved that these inequalities were
effective and they improved the performance of their branch-and-price
algorithm. However, the size of set $I_k$ is important and larger values
make the pricing algorithm more complicated. Since these cuts have already
been applied in a branch and price algorithm, we do not include these to
our study. We continue our investigation by looking at the projections
of the SPP-LRSP polytope.
%Their branch-cut-and-price algorithm
%using these inequalities is proved to be eff.

%\zaupdate{May include the part UPPER BOUNDING PROCEDURE using CLique Sepera
%tion function}

\subsection{Projection of the SPP-LRSP}
\label{pro-spp}
To eliminate the problem with the cuts (i.e.,dependency to specific columns
and thus passing dependent information for customers) that are valid
for the relaxation polytopes, we use projection. Here, projection is a
mapping of a polytope
to another based on a defined relation between variables.
%To be able to eliminate the dependency of a cut to specific columns, we define
%a projection of $\polyS$.
We let $v_{ji}=\sum_{p \in P_j}a_{ip}z_{p}$
for all $j \in J$ and $i \in I$ and project out the variables $z$. The projection
polytope $\polyPr$ defined based on this relation is
\begin{optprog}
  & \polyPr = \mbox{conv} & \{ t \in \{0,1\}^{|J|}, v \in \{0,1\}^{|J|\times|I|} | & \sum_{j \in J}v_{ji} & = & 1 & \forall i \in I,  \label{pro-sp1}\\
 &  &  & \sum_{i \in I} D_i v_{ji} -  L^F_{j} t_j & \leq & 0 & \forall j \in J,  \label{pro-sp2}\\
& & & & & & \qquad \ \   \}.  \nonumber
\end{optprog}
We can interpret $v_{ji}$ as the fraction of the demand of customer $i \in I$
assigned to facility $j \in J$.
%We can rewrite the constraints defining SPP-LRS using the variable transformation:
% \begin{optprog}
%           & v_{ji} & = &\sum_{p \in P_j}a_{ipj}Z_{jp}  & \forall i \in N,\forall j \in M  \label{transf}\\
%           & \sum_{j \in M} v_{ji} & = & 1 & \forall i \in N \label{trspp_con1} \\
%           & \sum_{i \in N} d_i v_{ji} & \leq & CapF_{j} T_{j} & \forall j \in M \label{trspp_con2} \\
%           & \sum_{j \in M} T_{j} & \geq & nREQ \label{tr-lb-open-fac}\\
%           & v_{ji} & \leq &  T_{j}  & \forall j \in M,i \in N \label{trrelation} \\
%           & v_{ji} &  \in & \{0,1\} & \forall j \in M, \forall i \in N \label{capbinary1} \\
%           & T_{j}  & \in & \{0,1\} & \forall  j \in M \label{capbinary2}
%\end{optprog}
Again, we investigate some relaxations of the polytope $\polyPr$ to derive
some cuts. Then, the obtained cuts can be projected back to the original variable
space $(t,z)$ and used for $\polyS$. We implemented two classes of the cuts
that will be presented next in this section: the lifted cover inequalities
and the generalized assignment polytope cuts. We extended our branch  and price
algorithm with these and presented some computational results in Section
\ref{comp:CoverGAP}.

\paragraph{Multiple Knapsack Polytope}
If we replace variable $t_j$ with 1 in inequalities
(\ref{pro-sp2}) and relax the equality restriction in
(\ref{pro-sp1}), these inequalities with the binary restriction for
variable $v$ define a multiple knapsack polytope:
\begin{optprog}
  & \polyMKn = \mbox{conv} & \{ v \in \{0,1\}^{|J|\times|I|} | & \sum_{j \in J}v_{ji} & \leq & 1 & \forall i \in I,  \label{pro-knp1}\\
 &  &  & \sum_{i \in I} D_i v_{ji} & \leq & L^F_{j} & \forall j \in J,  \label{pro-knp2}\\
& & & & & & \qquad \ \   \}.  \nonumber
\end{optprog}
Valid inequalities for the single knapsack polytope are also valid
for $\polyMKn$. The cover inequalities and (1-k) configuration
inequalities described in Section \ref{cuts-relax-spp} can be
applied here.

Let $S \subseteq I$ be a minimal cover for facility $j$ such that
$\sum_{i \in S}D_i> L^F_j$. Then, the following inequality is valid
for $\polyMKn$ and $\polyPr$.
\begin{equation}
\sum_{i \in S}v_{ji} \leq |S|-1. \label{inq-cover2}
\end{equation}
The inequalities (\ref{inq-cover2}) can be extended further to get
stronger inequalities using {\em lifting} which is a procedure to extend an
inequality to another possibly including nonzero coefficients for
some of the variables that have zero coefficient in the original
inequality. %$z_p$, for all $p \in P_j \backslash S_j$.
The extension of the cover inequalities and the lifting procedures
to get stronger inequalities are described in \citet{Nemhauser} and \citet{Gu}.
%Lifting is a procedure to extend the inequality (\ref{coverinq}) to also include
%variables $z_p$ for all $p \in P_j \backslash S_j$.
%As mentioned in the previous section, lifting procedure can be applied
%to inequality (\ref{inq-cover2}) to extend it to other variables in
%the formulation, thus to obtain possibly a stronger valid inequality.
To be more specific, after the lifting procedure, the inequality will be
as follows:
\begin{equation}
\sum_{i \in S}v_{ji} + \sum_{i \in I\backslash S}\alpha_{ji}v_{ji}
\leq \beta, \label{inq-cover2b}
\end{equation}
for some $\alpha$ and $\beta$ values that are calculated in a
lifting procedure. This inequality can be transformed to the following
for $\polyS$:
\begin{equation}
\sum_{i \in S}\sum_{p \in P_j}a_{ip}z_{p} + \sum_{i \in I\backslash
S}\sum_{p \in P_j}\alpha_{ji}a_{ip}z_{p} \leq \beta.
\label{inq-cover2c}
\end{equation}
An inequality in the form of (\ref{inq-cover2c}) can easily be added
to the RMP without destroying the pricing problem structure, i.e.,
our current pricing problem solution algorithms can be used. The
inequality (\ref{inq-cover2c}) has the information of whether a
customer visited by a pairing or not for all possible pairings (even
the ones that are not currently in RMP). The inequality may have
information about more than one customer, however this information
for a customer is independent of other customers.
% since the inequality depends on a set of customers not
%specific pairings, and all pairings (even the ones that are not
%currently in RMP) can have nonzero
%coefficient.  %if they include the customers in the set.
Let $\omega_S^j$ be the dual variable that corresponds to the
inequality (\ref{inq-cover2c}) for customer set $S$ and facility
$j$. To incorporate the information from this inequality into our
algorithm, in the pricing problem for facility $j$, we update the
cost of incoming arcs for customers in set $S$  with $\omega_S^j$,
and all other customers in $I\backslash S$ with $\omega_S^j
\alpha_{ji}$. These updates for single customers result an updated
total reduced cost for a column with a nonzero coefficient
in inequality (\ref{inq-cover2c}).

We implemented inequalities (\ref{inq-cover2c}) that are derived
from lifted cover inequalities in our branch and price algorithm.
MINTO has a built-in function to detect violated lifted cover
inequalities for general MIP instances.
%that is also called separation procedure.
This function implements a heuristic algorithm details of which is
described in \citet{Gu}. We export this function into our algorithm
%employ this function
to detect violated lifted cover inequalities for a given solution
information based on $(t,v)^j$ variable vector that corresponds to
some facility $j \in J$.

In our CUTGEN(RMP) procedure, for a given fractional solution
$(\bar{t},\bar{z}) \in {\mathbb{R}^{+}}^{|J|+\sum_{j \in J}|\bar{P}_{j}|}$
where $\bar{P}_{j}$ is the current set of pairings for facility
$j \in J$, we do the following.
\begin{itemize}
\item For each facility $j \in J$, if $\bar{t}_j = 1$, then we find
$\bar{v}^j \in \mathbb{R}^{|I|}$ using relation $\bar{v}_{ji} = \sum_{p \in
\bar{P}_j}a_{ip}\bar{z}_{p}$ for all $i \in I$.
\item For a given $\bar{v}^j$, we search for a cover inequality in the
form of (\ref{inq-cover2b}) using MINTO's function.
\item If the function returns a cut, then we rewrite
it in the form of (\ref{inq-cover2c}) and add to our RMP.
\end{itemize}
As mentioned earlier, computational results for the implementation
of this set of cuts are presented in Section \ref{comp:CoverGAP}.
%In addition, one can partition cover $S$ into two disjoint sets,
%$S_1$ and $S_2$, where $S_1 \cup S_2 = S$$S_1$

As mentioned in the knapsack polytope part in previous section,
(1-k) configuration inequalities are valid for single knapsack
polytope. Thus it is possible to derive (1-k) configuration
inequalities for the $\polyPr$ in terms of variable $v$. Similar to
cover inequalities, (1-k) configuration inequalities in terms of
variable $v$ can be added to the RMP without causing significant
changes in the pricing problem solution algorithms. Based on the
computational experiments for lifted cover inequalities which are a
special case of (1-k) configuration inequalities, we decided not to
implement (1-k) configuration inequalities in our study.

\paragraph{Generalized Assignment Polytope}
Similar to  the definition of $\polyMKn$,
% In constraint (\ref{pro-sp2}), by replacing multiplier $t_j$ with
%its upper bound 1, we obtain
a generalized assignment polytope, $\polyGAP$ can be defined
as follows:
 \begin{optprog}
  & \polyGAP = \mbox{conv}  \{ v \in \{0,1\}^{|J|\times|I|} | v \mbox{ satisfy (\ref{pro-sp1}) and } \sum_{i \in I} D_i v_{ji}  \leq  L^F_j \ \  \forall j \in J\}.  \nonumber
\end{optprog}

%\begin{eqnarray}
% X^{GAP}=& & \{v \in \mathbb{Z}^{|M|X|N|}_+ | v  \mbox{ satisfy const. (\ref{trspp_con1})},\\
%& & \sum_{i \in N} d_i v_{ji} \leq  CapF_{j} \ \ \forall j \in M, \nonumber \\
%& &  v_{ji}\leq 1 \ \ \forall j \in M \}\nonumber
%\end{eqnarray}
The generalized assignment problem (GAP) is a generalization of the
multiple knapsack problem. The generalization is that the GAP may be
formulated with different weights for each item (coefficients in
\ref{pro-knp2}, i.e.,$D_i$) for each knapsack.
Thus, valid inequalities for the knapsack polytopes are also valid for
$\polyGAP$.
%Since $X^{GAP}$ is a relaxation of the convex hull of the SPP-LRS problem, the valid inequalities for
%$X^{GAP}$ are also valid for the SPP-LRS polytope.
Besides the inequalities for the knapsack polytope, there are only a
few papers investigating the facets and valid inequalities of GAP.
Some of these are \citet{Gottlieb90}, \citet{Farias}.
Both papers relax constraint (\ref{pro-sp1}) and look at the relaxed polytope
including $\leq $ part of inequalities (\ref{pro-sp1}). In this
relaxed GAP, \citet{Farias} added the condition that at most one of
$v_{ji}$ for all $j \in J$ can be positive for all $i \in I$. The
following theorem from \citet{Farias} presents a valid inequality
for GAP. We rewrite the theorem based on our notation and variable
definition.
\begin{thm}
\label{GAP-cut}
[\citet{Farias}] Let $S\subseteq I$ and $k\in J$. Suppose
that $\sum_{i \in S}d_i > L^F_k$ and $\sum_{i \in S-\{r\}}d_i < L^F_k$ for
some $r \in S$. Then,
\begin{optprog}
& \sum_{i \in S}D_i v_{ki} +  \sum_{i \in S} \left(\mbox{max}\{0,L^F_k- \sum_{a \in S-\{i\}}D_a \}\right)\sum_{j \in J-\{k\}}v_{ji}  & \leq &  L^F_k \label{gap-valid}
\end{optprog}
 defines a facet of relaxed GAP polytope.
\end{thm}
The interpretation of (\ref{gap-valid}) is simple: if the total
demand of the customers in set $S$ exceeds the capacity of facility
$k$, then there must be some customer $r \in S$ that must be served
by another facility and not by $k$. If the second part of the left
hand side of (\ref{gap-valid}) is zero, it is equal to the
facility capacity constraint for $k$ that is one of the trival
inequalities. However, in a non integer solution, if the slack
variable for the facility capacity constraint for $k$ is zero (i.e.,
the capacity is fully used) and facility $k$ has positive flow to
all customers in set $S$, then the inequality (\ref{gap-valid}) may
likely be violated because of positive flow from other facilities to
some customer $r \in S$ with $L^F_k- \sum_{a \in S-\{r\}}d_a > 0$.
We can rewrite inequality (\ref{gap-valid}) in terms of pairing
variables using the relation between $v$ and $z$:
\begin{optprog}
& \sum_{i \in S}d_i \sum_{p \in P_k}a_{ip}z_{p} +  \sum_{i \in S} \left(\mbox{max}\{0,L^F_k- \sum_{a \in S-\{i\}}D_a \}\right)\sum_{j \in J-\{k\}}\sum_{p \in P_j}a_{ip}z_{p}  & \leq &  L^F_k \label{gap-valid2}
\end{optprog}
%\begin{eqnarray}
%& &  \sum_{i \in S}\sum_{p \in P_t}d_i a_{itp}Z_{tp} +  \nonumber\\
%& &  \sum_{i \in S} \left(max\{0,CapF_t-
%\sum_{x \in S-\{i\}}d_x \}\right)\sum_{j \in M-\{t\}}\sum_{p \in P_j}d_i a_{ijp}Z_{jp}  \nonumber \\
%& & \hspace{3.5in} \leq  CapF_t \label{gap-valid}
%\end{eqnarray}

%\ref{GAP-cut} is a valid inequality for our problem.

Since the inequality (\ref{gap-valid2}) carries independent
information of presence of a customer in a pairing for each customer
in set $S$, it can be added to the RMP without destroying the
structure of the pricing problem. We implemented this class of
inequalities which we refer as GAP inequalities.
We extended our branch and price algorithm with (\ref{gap-valid2})
and evaluate the overall performance. Let $\gamma_S^k$ be the dual variable
corresponding to the inequality (\ref{gap-valid2}) for customer set
$S$ and facility $k$. To incorporate this inequality into our
algorithm, in the pricing problem for facility $k$, we update the
incoming arcs for customers in set $S$  with $D_i \gamma_S^k$, and
in the pricing problem for facility $j \neq k$, we update the
incoming arcs for customers in set $S$ with
$(\mbox{max}\{0,L^F_t-\sum_{a \in S-\{i\}}D_x\})\gamma_S^k$.

\citet{Farias} present the following separation heuristic that is
described based on our notation for inequality (\ref{gap-valid}):
%The steps of the heuristic are as follows:
\begin{center}
\fbox{\begin{minipage}{5.6in}
{\bf Heur-Sep-GAP [\citet{Farias}]}
\begin{enumerate}
\item[0.] Let $\bar{v}$ be a fractional solution to the LP relaxation of $\polyGAP$.
\item[1.] For each facility $k \in J$:
      \begin{itemize}
      \item[i.] If $\sum_{i \in I}D_i\bar{v}_{ki} = L^F_k$, then proceed.
      Otherwise go to step 1 and pick the next facility.
      \item[ii.] Let $S = \{i \in I | \bar{v}_{ki} >0 \}$ and  $D(S)= \sum_{i \in S}D_i$.\\
      If $D(S) > L^F_k$, then for each $c \in S$, check the following conditions:
      \begin{itemize}
      \item Is $c$ also served by another facility?
      \begin{optprog}
      & 1 > v_{kc} > 0, \mbox{ and }\sum_{j \in J\backslash \{k\}} v_{jc} & > & 0 \label{her-sep-con1}
      \end{optprog}
      \item When we remove $c$ from assignment set of facility $k$,
      does full assignment of the rest of the customers to $k$ becomes feasible?
      \begin{optprog}
      & D(S \backslash \{c\}) & < & L^F_k \label{her-sep-con2}
      \end{optprog}
      \end{itemize}
%      \begin{optprog}
%      & 1 > v_{ka} > 0, \mbox{ thus } \sum_{j \in J\backslash \{k\}} v_{ja} & > & 0 \label{her-sep-con1} \\
%       & D(S \backslash \{a\}) & < & L^F_k \label{her-sep-con2}
%      \end{optprog}
          %\begin{itemize}
         %\item $ 1 > v_{ka} > 0$, thus $\sum_{j \in J\backslash \{k\}} v_{ja} > 0$.
         %\item $D(S \backslash \{a\}) < L^F_k$.
         %\end{itemize}
     \item[iii.]
     %If there is at least one $c \in S$ satisfying (a) and (b),
    %   then (\ref{gap-valid}) is violated for facility $k$ and customer set $S$.
     If there is at least one $c \in S$ satisfying both (\ref{her-sep-con1}) and
     (\ref{her-sep-con2}), then (\ref{gap-valid}) is violated for facility $k$ and
       customer set $S$. Return a GAP cut for facility $k$.
     \end{itemize}
\end{enumerate}
\end{minipage}}
\end{center}

In our CUTGEN procedure, similar to the lifted cover inequalities,
for a given fractional solution
$(\bar{t},\bar{z}) \in {\mathbb{R}^{+}}^{|J|+\sum_{j \in J}|\bar{P}_{j}|}$
where $\bar{P}_{j}$ is the current set of pairings for facility $j \in J$, we
find $(\bar{t},\bar{v}) \in \mathbb{R}^{|J|+|J|\times|I|}$ using relation
$\bar{v}_{ji} = \sum_{p \in \bar{P}_j}a_{ip}\bar{z}_{p}$ for all $i \in I$ and
$j \in J$. Then, we call
%In our implementation, we also use
Heur-Sep-GAP to separate cuts in the form of inequality (\ref{gap-valid}).

We updated Heur-Sep-GAP slightly in our implementation.
In step 1(i), we also require that the facility variable is 1, i.e.,
$t_k = 1$. Then, we rewrite a generated inequality in the form of
(\ref{gap-valid2}). In the LRSP instances, most of the time Heur-Sep-GAP
fails to find violated cuts. The reason for this failure is condition
(\ref{her-sep-con2}). Generally, condition (\ref{her-sep-con2}) is not
satisfied even though the facility capacity constraint is binding and
(\ref{her-sep-con1}) is true. In LRSP, we are dealing with the assignment
of pairings (i.e., sets of customers to the facilities), therefore,
fractional pairing variables lead to multiple customers (rather than a single
customer) that are served by
more than one facility. In a fractional solution, for some facility,
even though the capacity is fully utilized, since more than one customer
is served by a fractional amount from that facility, it is not enough to
remove the assignment of a single customer to this facility to satisfy that
the total demand of the remaining customers is less than the facility capacity,
(i.e.,condition (\ref{her-sep-con2})). To increase the number of times a cut is
generated, we updated determination of the set $S$ in the step 1 (ii) of the
Heur-Sep-GAP. We may shrink this set $S$ whenever it is necessary.

In our algorithm, first we process step 1 (ii) and (iii) as in the original
Heur-Sep-GAP procedure by setting $S = \{i \in I | v_{ki} >0 \}$.
If we cannot find a violation, we reduce the size of set $S$
by removing one customer at a time until a violation is found
or $D(S) \leq L^F_k$.
%We reprocess step 1 (ii) in our implementation.
%First, we set $S = \{i \in I | v_{ki} >0 \}$, then execute
%the step. If we can not find a violation,
%We update $S$ by removing one customer at a time until we get a violation
%with the updated set or $D(S) \leq L^F_k$.
We remove customers in the order of non increasing demand.
The steps of the updated algorithm can be listed as follows:
\begin{center}
\fbox{\begin{minipage}{5.6in}
{\bf Heur-Sep-GAP-2}
\begin{enumerate}
\item[0.] Let $(\bar{t},\bar{v})$ be a fractional solution.
\item[1.] For each facility $k \in J$:
    \begin{itemize}
      \item[i.] If $\bar{t}_k = 1$ and $\sum_{i \in I}D_i\bar{v}_{ki} = L^F_k$, then
      proceed. Otherwise return to step 1.
      \item[ii.] Let $S = \{i \in I | \bar{v}_{ki} >0 \}$. If $D(S) > L^F_k$,
      then for each $c \in S$, check the conditions (\ref{her-sep-con1}) and
     (\ref{her-sep-con2}).
     \item[iii.] If there is at least one $c \in S$ satisfying (\ref{her-sep-con1}) and
     (\ref{her-sep-con2}), then return  a cut for facility $k$ and
       set $S$ and return to step 1.
     \item[iv.] Sort the customers in $S$, $i_1<i_2<..<i_{|S|}$ such that
     $D_{i_p} \geq D_{i_q}$ for all $p, q = 1,..,|S|$ and $i_p<i_q$.
     \item[v.] For $e = 1$ to $|S|$:
       \begin{itemize}
        \item[a.] Let $S = S - \{i_e\}$. If $D(S) \leq L^F_k$, then return to step 1.
        \item [b.] If there is at least one $c \in S$ satisfying
        (\ref{her-sep-con1}) and (\ref{her-sep-con2}), then return
        a cut for facility $k$ and set $S$ and
        return to step 1.
        \end{itemize}
     \end{itemize}
\end{enumerate}
\end{minipage}}
\end{center}
As mentioned earlier, computational results for the implementation
of this set of cuts are presented in Section \ref{comp:CoverGAP}.
%Should I briefly mention implementation results?}

\paragraph{Capacitated Facility Location Polytope}
In $\polyPr$, if we relax the integrality of variable $v$ and
set bounds, we get a capacitated facility location polytope, $\polyFL$. To be
more specific:
\begin{optprog}
  & \polyFL = \mbox{conv} & \{ t \in \{0,1\}^{|J|}, v \in \mathbb{R}^{|J| \times |I|} | t, v \mbox{ satisfy } (\ref{pro-sp1}), (\ref{pro-sp2}), \mbox{ and }v_{ji} \in [0,1] \forall j \in J, i \in I\}.  \nonumber
\end{optprog}
Therefore, $\polyFL$ is a relaxation of $\polyPr$, and if the valid inequalities
for $\polyFL$ are projected to $(t,z)$ space, they are valid for $\polyS$.
Indeed, $\polyPr$, the capacitated facility location polytope with binary
restriction on variable $v$, is called the single source capacitated facility location
polytope. In the problem, each customer has to be served by a single facility.
We do not aware of any studies that investigate specifically the polyhedral structure
of single source capacitated facility location problem, thus we review the studies
for $\polyFL$ which are a few.  \citet{Aardal}, \citet{Leung} and
\citet{Padberg} investigated valid inequalities and gave conditions for these
inequalities to be facets for $\polyFL$. We did not implement any of these
inequalities, here we give a brief list of common ones.

One set of valid inequalities is the {\em knapsack cover inequality} in terms
of $t$ (facility) variables. \citet{Aardal} define "cover" that is a subset of
facilities. Using our notation, $M$ is a cover if
$\sum_{j \in M}L^F_j>\sum_{j \in J}L^F_j-\sum_{i \in I}D_i$. Based
on this definition, they state the following theorem:
\begin{thm}
\label{thm-kpcover}
[\citet{Aardal}] If $M$ is a minimal cover with respect to $J$ and $I$,
$L^F_{\mbox{min}} = \mbox{min}_{j \in M}\{L^F_j\}$,
and $\sum_{j \in J \backslash M}L^F_j + L^F_{\mbox{min}} > \sum_{i \in I}D_i$, then the
knapsack cover inequality:
\begin{equation}
\label{Tcover}
\sum_{j \in M}t_j \geq 1
\end{equation}
defines a facet of $\polyFL \cap \{t \in \{0,1\}^{|J|}:t_j=1 \qquad \forall j\in J \backslash M\}$.
\end{thm}

This theorem implies that at least one facility from set $M$ must be open
because the total capacities of facilities in the set $J \backslash M$ is not
enough to cover the total demand of customers. Since \ref{Tcover}
is valid for $\polyFL$, then it is also valid for $\polyPr$ and $\polyS$.
Since the inequality defined in Theorem (\ref{thm-kpcover}) only includes $t$ variables,
in other words it is a z-free inequality, we can use this inequality
without affecting the structure of the pricing problem.

We did not implement this inequality in our study because of the effect of
inequality (\ref{lb-open-fac}) in the RMP. In our LRSP instances used for
computational results, the facility capacity values are equal. Based on this, when all
possible (\ref{Tcover}) inequalities are generated, summed  and combined
into one inequality, we get $\sum_{j \in J}t_j \geq m$ where $m \in \mathbb{R}^+$.
With some algebra, we can easily show that the minimum number of required
facilities, $R$ (in (\ref{lb-open-fac})) is greater than $m$, if $1 < R < |J|$
(which is the case in our instances). Thus, the inequality (\ref{lb-open-fac})
forcing that the total number of facilities must be greater than or equal to
$R$ is stronger than the inequality obtained by summing all possible inequalities
(\ref{Tcover}). In addition, we do not add the inequalities (\ref{Tcover})
to the SPP-LRSP, since generally inequality (\ref{lb-open-fac}) results
integer values for location variables.

Another class of valid inequality derived for $\polyFL$ is {\em flow cover inequality}.
\citet{Aardal} and \citet{Padberg} define $M$ as flow cover if $M$ is a subset of
facilities and can strictly cover the total demand of customers, i.e.,
$\sum_{j \in M}L^F_j = \sum_{i \in I}D_i + \lambda$, for some $\lambda > 0$.
\citet{Aardal} and \citet{Padberg} present the following theorem
(based on the notation in this section) that describes the flow cover
inequalities:
\begin{thm}
\label{thm-fwcover}
[\citet{Aardal}]
Let $M \subset J $ be a flow cover with respect to $J$ and $I$.
%i.e.,$\sum_{j \in M}L^F_j = \sum_{i \in I}D_i + \lambda$, for some $\lambda > 0$
If $\mbox{max}_{j \in M}\{L^F_j\} > \lambda$ and
$\sum_{j \in J}L^F_j > \sum_{i \in I}D_i + L^F_r$ $\forall r \in M$, then
the flow cover inequality
\begin{equation}
\label{FcoverAardal}
\sum_{j \in M}\sum_{i \in I}D_i v_{ji}+\sum_{j \in M}(L^F_j-\lambda)^+(1-t_j) \leq \sum_{i \in I}D_i
\end{equation}
defines a facet of $\polyFL$.
\end{thm}
The inequality (\ref{FcoverAardal}) defines a minimum set of facilities (i.e., $M$)
to cover the total demand. If all customers are assigned to the facilities in $M$, the
excess capacity $\lambda$ leads at least one facility in set $M$ to be fractional in an LP
solution. The aim of the inequality (\ref{FcoverAardal}) is to eliminate the effect of
excess capacity. The second part in the left hand side of the inequality
forces all facilities in set $M$ take value close to 1 to satisfy the inequality.
In inequality (\ref{FcoverAardal}), if we replace $v_{ji}$ variable with
$\sum_{p \in P_j}a_{ip}z_{p}$, we have the following inequality in $(t,z)$
space.
\begin{equation}
\label{FcoverZvar}
\sum_{j \in M}\sum_{p \in P_j}\sum_{i \in I}D_i a_{ip}z_p + \sum_{j \in M}(L^F_j-\lambda)^+(1-t_j) \leq \sum_{i \in I}D_i
\end{equation}
Similar to the other inequalities defined in this section, we can incorporate the inequality
(\ref{FcoverZvar}) into a branch and price algorithm. However, we did not implement this
inequality in our study, because in most of our LP relaxation solutions, the facilities
serving the customers are fully open.
%Let $\theta_M$ be the dual variable assigned to the inequality when it is added to the master
%problem. The reduced cost of pairings for facilities in the flow cover will be updated with
%$\theta_M$, and thus in the pricing problem, the costs of the incoming arcs to customer nodes
%are updated with $D_i\theta_M$. We did not implement this inequality in our algorithm.

%\begin{eqnarray}
%  \label{new-red-cost-fc}
%  \hat{C}_{p}=C_{p}-\sum_{i\in I}a_{ip}\pi_{i} + \sum_{i \in I}a_{ip}D_{i}(\mu_{j}+\theta_J)+\sum_{i \in I}a_{ip}\sigma_{ji}
%\end{eqnarray}
%In the pricing problem of facility $j \in J$, we have to adjust the cost of the incoming
%arcs to customer nodes with  $(\theta_J)$.
%The cost of an arc $(i,k)$ should be equal to $c_{ik}-\pi_{k}+d_k \cdot (\mu_j+\theta_J)+\sigma_{jk}$.
%The pricing problems for facilities that are not in set $J$ do not change.
\begin{comment}
The flow cover inequality (\ref{FcoverZvar}) uses a flow cover with respect to
the set of all customers, $I$. \citet{Aardal}
introduce {\em effective capacity inequality} that uses a flow cover with respect to a
subset of customers. They prove that the effective capacity inequality is valid,
and defines a facet for $\polyFL$ under some  conditions. Let $M \subset J$ be a flow cover
with respect to $K \subseteq I$, then $\sum_{j \in M}L^{F'}_{j}=\sum_{i \in K}D_i + \lambda$
for some $\lambda>0$  and $L^{F'}_j=\mbox{min}\{L^F_j,\sum_{i \in K_j}D_i\}$ for some sets
$K_j\subseteq K$ for all $j \in M$.
\begin{prop}
\label{prop-effwcover}
[\citet{Aardal}] Let $M \subset J$ be a flow cover with respect to $M$ and $K$.
If $\mbox{max}_{j \in M}\{L^{F'}_j \} > \lambda$ then the following effective flow cover inequality
\begin{equation}
\label{EFFcover}
\sum_{j \in M}\sum_{i \in K_j}D_i v_{ji}+\sum_{j \in M}(L^{F'}_j-\lambda)^+(1-t_j) \leq \sum_{i \in K}D_i
\end{equation}
is valid for $\polyFL$.
\end{prop}
The inequality (\ref{EFFcover}) can also be incorporated into the branch and price
algorithm in a similar way with inequality (\ref{FcoverZvar}). When the inequality is added
to the master problem, in the pricing problem of facility in the flow cover, only the costs of the
incoming arcs to customers in set $K_j$ are updated with the dual variable. Again, we did not
implement effective capacity inequalities in this study.
\end{comment}
%\zaupdate{flow cover inequalities implemented in (\citet{VanRoy}) ??}



\section{Disjunctive Cuts for LRSP}
\label{disjcuts-sec}
In this section, we discuss the generation of a set of cuts derived based on split
disjunctions in the context of branch and price algorithm.
We begin with a review of the methodology of disjunctive cuts. The generation of 
disjunctive cuts and the idea of disjunctive programming
were introduced by \citet{Balas79}. Disjunctive programming is finding the optimal
value of an objective function over a finite union of polyhedra constructed based
on conjunctive statements (connected with logical operator "AND") and disjunctive
statements (connected with logical operator "OR"). Here, we focus on cuts based on
disjunctive statements. To explain the methodology for an integer
model, suppose we have an integer program, P:
\begin{optprog}
   Minimize & \objective{c^Tx} \nonumber\\
(P) \ \  subject to & Ax & \geq & d, \label{dp-1}\\
                     & x & \geq & 0,\mbox{ and integer}  \label{dp-2}
\end{optprog}
where $A \in \mathbb{Q}^{m\times n}$ is constraint matrix, $c \in \mathbb{Q}^{n}$,
$d \in \mathbb{Q}^{m}$ are rational vectors and $x \in \mathbb{Z}^n$ is a vector
of decision variables. In addition, suppose we know that any feasible integer solution
$x$ for P, also satisfies at least one of the following $s$ inequalities based
on some disjunctive statements:
\begin{optprog}
%\sublabon{optprog}
& B^1x & \geq & b^1, &  \label{dissys1}\\
& B^2x & \geq & b^2, &  \label{dissys2}\\
& ... \nonumber\\
& B^sx & \geq & b^s, &  \label{dissys3}
%\sublaboff{optprog}
\end{optprog}
where $B^k \in \mathbb{Q}^{m^k\times n}$ is constraint
matrix and $b^k \in \mathbb{Q}^{m^k}$ is right hand side vector
for some integer $m^k > 0$, $\forall k = 1,..,s$.
Using these $s$ disjunctive set of inequalities, a relaxation of P, P$_{\mbox{DR}}$
can be defined:
\begin{optprog*}
   Minimize & \objective{c^Tx} \nonumber\\
(P$_{\mbox{DR}}$) \ \  subject to & Ax & \geq & d, \\
 & {\lor}_{k = 1}^{s}B^kx & \geq & b^k, \\
 & x & \geq & 0,
% & x & \in & {\mathbb{R}^+}^n,
\end{optprog*}
where $\vee$ denotes ``or'' and is used to represent the disjunctive system.
Using each disjunction, we can define $s$ polyhedron.
\begin{optprog*}
& \mbox{P}^k = \{x \in  {\mathbb{R}^{+}}^n | Ax \geq d, B^kx  \geq  b^k\} & \forall k = 1,..,s
\end{optprog*}
Based on the definitions of P$^k$ and P$_{\mbox{DR}}$, P$_{\mbox{DR}}$ can be represented
by the union of these $s$ polyhedron, i.e.,
P$_{\mbox{DR}} = \cup_{k = 1}^s\mbox{P}^k$. Since P$_{\mbox{DR}}$ is a relaxation
of P, any valid inequality for P$_{\mbox{DR}}$ is also valid for P and
P$\subseteq$P$_{\mbox{DR}} \subseteq$P$_{\mbox{LP}}$, where P$_{\mbox{LP}}$ is the LP relaxation
of P. The aim of the disjunctive cut generation methodology is to drive an inequality,
$\alpha x \geq \beta$ that is valid for P$_{\mbox{DR}}$ and is violated by some points of
P$_{\mbox{LP}}$.

\citet{Balas79} show that based on the assumption that $\mbox{P}^k$ is non empty for all
$k = 1,..,s$, $\alpha x \geq \beta$ is valid for P$_{\mbox{DR}}$ if and only if there exist
multipliers $u^k \in \mathbb{Q}^{m}$ and $v^k \in \mathbb{Q}^{m^k}$ such that
\begin{optprog}
& u^kA & + & v^kB^k & \leq & \alpha, & \forall k = 1,..,s \label{dcutgen1}\\
& u^kb & + & v^kb^k &  =   & \beta,  & \forall k = 1,..,s \label{dcutgen2} \\
& & &  u^k, v^k &\geq & 0.  & \forall k = 1,..,s \label{dcutgen3}
%& \alpha & \geq & u^kA & + & v^kB^k, & \forall k = 1,..,s \\
%& \beta & = & u^kb & + & v^kb^k, & \forall k = 1,..,s  \\
%&  u^k, v^k &\geq & 0.  & & & \forall k = 1,..,s
\end{optprog}
Let $\bar{x}$ be a solution to the LP relaxation of P. To find a cut
$\alpha x \geq \beta$ that is maximally violated by $\bar{x}$ and valid for
P$_{\mbox{DR}}$ (separation problem), we need to solve the following optimization
model, which is called {\em cut generating LP} (CGLP).
\begin{optprog}
(CGLP) Minimize & \objective{\alpha \bar{x} -  \beta} \label{cglp1}\\
& u^kA & + & v^kB^k & \leq & \alpha, & \forall k = 1,..,s \label{cglp2}\\
& u^kb & + & v^kb^k &  =   & \beta,  & \forall k = 1,..,s  \\
& \sum_{i=1}^n u^k_i & + & \sum_{i=1}^n v^k_i & = & 1, \label{norm1}\\
& & &  u^k, v^k &\geq & 0.  & \forall k = 1,..,s  \label{cglpm}
%& \alpha & \geq & u^kA & + & v^kB^k, & \forall k = 1,..,s \\
%& \beta & = & u^kb & + & v^kb^k, & \forall k = 1,..,s  \\
%&  u^k, v^k &\geq & 0.  & & & \forall k = 1,..,s
\end{optprog}
Constraint (\ref{norm1}) is added to normalize the cut to be
generated. The purpose of the normalization constraint is to
truncate the cone defined by constraints (\ref{dcutgen1}) -
(\ref{dcutgen3}) to a polyhedron. The normalization removes the
unboundedness of the model and transforms extreme rays to
extreme points. Here, we use one of the common normalization
constraints, different versions of normalization
constraints are described in the literature. Some examples are
$\alpha -$normalization (i.e.,$\sum_{i=1,..,n}|\alpha_i| = 1$) and
$\beta -$normalization (i.e.,$|\beta| = 1$) (\citet{Balas96}).

As noted in \citet{Caprara03} (based on the results in \citet{Balas93} and
\citet{Balas96})), for a fixed disjunctive system, i.e.,
(\ref{dissys1}) - (\ref{dissys3}), separation of disjunctive cuts
can be done in polynomial time.

In general, to derive disjunctive cuts, it is common to use disjunctive
systems including two linear inequalities (i.e.,$s = 2$).
%One of the most common and successful general purpose cutting planes for
%0-1 mixed integer programming models is
One of the most common and successful set of disjunctive cuts is referred to as
{\em lift-and-project cuts} (see \citet{Balas93} and \citet{Balas96}) that are
derived from bounds of single binary variables. The disjunctions used to generate
lift-and-projects cuts are in the form of $(x_i \leq 0) \vee (x_i \geq 1)$ for some
binary variable $x_i$. Lift-and-project cuts are one of the widely used
general purpose cutting planes for 0-1  mixed integer programming models.
For models with a polynomial number of constraints, separation of these cuts
can be done in polynomial time by solving the cut generating LP. In addition to this
advantage, the problem of having to solve large cut generating LPs for many iterations
in a branch-and-cut based algorithm was eliminated with results presented in
\citet{Balas03}. They showed that instead of explicitly formulating and solving CGLP,
the simplex tableau of the current LP solution can be used to derive this type of cuts.


If the disjunctions used to derive cuts include more than one variable and
are in the form of $(\pi x \leq \pi_0) \vee (\pi x \geq \pi_0 + 1)$ for some
integer vector $(\pi,\pi_0)\in \mathbb{Z}^{n+1}$, they are called {\em split disjunctions}.
And the cuts generated are referred to as {\em split cuts} (\citet{Cook90} and \citet{Balas08}).
\citet{Caprara03} proved that the separation problem for split cuts is strongly NP-Complete.
Many of the general purpose cutting planes for mixed integer programming models
can be included in the class of split cuts. Lift-and-project cuts are simpler forms of
split cuts. In addition,  mixed integer cuts of Gomory (\citet{Gomory63}) and mixed integer
rounding inequalities (\citet{Nemhauser90}) are shown to be in the class of split cuts
(\citet{Cornuejols01} and \citet{Cornuejols06}).

In addition to these studies that derive cutting planes based on some disjunctions,
\citet{Letchford01} proved that  many facet defining inequalities for some well-known
combinatorial optimization problems investigated before are actually disjunctive cuts
derived from certain simple disjunctions.
%, many facet defining inequalities for some of the combinatorial
%optimization problems investigated before are proved to be disjunctive cuts drived from certain
%simple disjunctions (\citet{Letchford01}).
They used these results to prove that the separation
problems for these facet inducing inequalities are solvable in polynomial time. They provided
results for some facets of the clique partitioning problem, the maximum cut problem,
the acyclic sub digraph problem, the linear ordering problem, the asymmetric traveling
salesman problem and the set covering problem.
%2-chorded odd cycle inequalities (\citet{Grotschel90}) and odd wheel inequalities
%(\citet{Chopra93}) for clique partitioning polytope,

As mentioned before, many general purpose cutting planes (such as Gomory cuts, mixed
integer rounding cuts) for mixed integer programming models can be  derived based on
disjunctions and can be included in the class of disjunctive cuts. Although
these general purpose cutting planes are very popular in a branch and cut based
algorithm, we haven't seen a study in the literature that is utilizing disjunctive cuts
in a branch and price algorithm before, except in a recent
presentation describing the application of disjunctive cutting planes in a branch and price
algorithm developed for capacitated VRP (\citet{Ropke08}). Next, we describe some
implementations of disjunctive cutting planes designed for our branch and price algorithm
for LRSP. Since we are applying a column generation algorithm, we avoided
using formulation SPP-LRSP to generate disjunctive cuts. Instead we used some projections
and relaxations of SPP-LRPS and G-LRSP. In addition,  the structure of the generated
cutting plane is crucial to be able to include it to the column generation algorithm.
We designed four implementations generating disjunctive cuts based on some formulations
of LRSP. The main steps of the implementations  are similar and can be outlined as follows:
%\begin{minipage}[h]{0.9\linewidth}

\begin{center}
\fbox{\begin{minipage}{5.6in}
\begin{enumerate}
\item[]  DISJCUTGEN:
\item[0.]  Let $\bar{y}$ be a given fractional solution to the RMP (at any node).
\item[1.]  Choose a model (representing a relaxation of the LRSP or its projection based
on some formulation (SPP-LRSP or G-LRSP)) to derive cuts. Let P be the selected model (as defined  in this section).
\item[2.] If P is from a projection, let $\bar{x}$ be the solution to P corresponding to
$\bar{y}$.
\item[3.] Select two disjunctive inequalities violated by $\bar{x}$ but
satisfied by all integer feasible solutions (we list possible disjunctions to be
considered for each implementation).

Let $B^1x \geq b^1 \lor B^2x \geq b^2$
be the selected disjunction.
\item[4.] Construct CGLP as defined by constraints (\ref{cglp2}) - (\ref{cglpm})
and objective (\ref{cglp1}).
\item[5.] Solve CGLP. If the objective function value is negative, then
$\alpha x \geq \beta$ is a valid cut violated by $\bar{x}$. Else STOP with no cut
generated.
\item[6.] If P is from a projection, project the generated cut into the original space,
let $\alpha' y \geq \beta'$ be the obtained cut.
\end{enumerate}
\end{minipage}}
\end{center}
%We tried implementation

In each of four implementations, step 1 is different. The procedures start with different
models and generate cuts based on the selected model. The details of each implementation
are explained next.


\subsection{Implementation 1}
\label{sec:djc-imp1}
In this implementation, we use the projection of $\polyS$, $\polyPr$ described in
Section \ref{pro-spp}.
Let $(\bar{t}, \bar{z}) \in {\mathbb{R}^{+}}^{|J|+\sum_{j \in J}|\bar{P}_{j}|}$ be a
fractional solution to the RMP (at any node of the tree, after column generation
algorithm is complete) where $\bar{P}_{j}$ is the current set of pairings for facility 
$j \in J$. Then we can find $(\bar{t}, \bar{v}) \in {\mathbb{R}^{+}}^{|J|+|J|\times|I|}$
with relation $\bar{v}_{ji}  =  \sum_{p \in \bar{P}_j}a_{ip}\bar{z}_{p}$
for all $i \in I$, $j \in J$.

We check the following disjunctions to find two disjunctive inequalities violated by the
fractional solution $(\bar{t}, \bar{v})$.
\begin{enumerate}
\item Flow from a facility to a customer is either 1 or 0:
%$v_{ji} \leq 0 \vee v_{ji} \geq 1$ for all $j \in J$, $i \in I$.
\begin{optprog}
& v_{ji} \leq 0  & \vee & v_{ji} \geq 1 & \forall j \in J, i \in I. \label{imp1-disj1}
\end{optprog}
\item Total number of customers assigned to a facility is an integer:
%$\sum_{i \in I}v_{ji} \leq g^j \vee \sum_{i \in I}v_{ji} \geq g_j + 1$
%for some integer $g^j$ and for all  $j \in J$.
\begin{optprog}
& \sum_{i \in I}v_{ji} \leq g_j & \vee & \sum_{i \in I}v_{ji} \geq g_j + 1 & \forall j \in J, \label{imp1-disj2}
\end{optprog}
for some integer $g_j$ and for all  $j \in J$.
\end{enumerate}

%\zaupdate{How we choose the disjunctions, close to 0.5,}
%\zaupdate{ parameters MAXNUM\_ROUNDS\_D1, MAXNUM\_ROUNDS\_D2}

After deciding on the form (\ref{imp1-disj1}) or (\ref{imp1-disj2}) of the
violated disjunctive inequalities, we write the CGLP
((\ref{cglp1}) - (\ref{cglpm})) with $s = 2$ based on model $\polyPr$. As we noted in
the description of the branch and price
algorithm, we use AUG-SPP-LRSP in the algorithm, thus $\polyPr$ is augmented with
projections of inequalities (\ref{lb-open-fac}) and (\ref{relation}). The LP relaxation of
the projected model used for the CGLP is as follows:
\begin{optprog}
{\bf [$M_1$]} & \sum_{j \in J}v_{ji} & = & 1 & \forall i \in I,  \label{m-pro-sp1}\\
& \sum_{i \in I} D_i v_{ji} -  L^F_{j} t_j & \leq & 0 & \forall j \in J,  \label{m-pro-sp2}\\
& \sum_{j \in J}t_j & \geq & R, \label{m-pro-sp3}\\
& v_{ji} - t_j & \leq & 0  & \forall j \in J, \forall i \in I, \label{m-pro-sp4}\\
& 0 \leq t_j & \leq & 1 & \forall j \in J,\\
& 0 \leq v_{ji} & \leq & 1 & \forall j \in J, \forall i \in I.
\end{optprog}
We investigate a disjunctive cut that is maximally violated by $(\bar{t}, \bar{v})$ 
in the following form:
\begin{equation}
\label{imp1-dc}
\sum_{j \in J}\alpha^t_{j}t_j + \sum_{j \in J}\sum_{i \in I}\alpha^v_{ji}v_{ji} \geq \beta.
\end{equation}
The objective function of the CGLP is to minimize
$\sum_{j \in J}\alpha^t_{j}t_j + \sum_{j \in J}\sum_{i \in I}\alpha^v_{ji}v_{ji} - \beta$.
The solution to the CGLP provides coefficients
$(\alpha, \beta) \in {\mathbb{R}}^{|J|+|J|\times |I|+1}$, and we can rewrite a
cut in the form of (\ref{imp1-dc}) in the $(t,z)$ variable space:
\begin{equation}
\sum_{j \in J}\alpha^t_{j}t_j + \sum_{j \in J}\sum_{p \in P_j}\sum_{i \in I}\alpha^v_{ji}a_{ip}z_{p} \geq \beta.  \label{disj_cut_z}
\end{equation}
A cut in the form of (\ref{disj_cut_z}) can easily be added to the column generation
algorithm by adding cost $\alpha^v_{ji}\omega$ to the cost of the links to
customer $i$ for all $i \in I$ in the pricing problem for facility $j \in J$ where
$\omega$ is the dual variable that corresponds to (\ref{disj_cut_z}).

%Computational results for this implementation are presented in Section
%\ref{sec:disjcutres} and a discussion of the results.
%\zaupdate{mention results of implementation 1?} % Not working? why??}
%\zaupdate{outline of the algorithm}

\subsection{Implementation 2}
In the second implementation of DISJCUTGEN, in step 1, we choose a graph based 
formulation which is a projection of a relaxation of G-LRSP.
%Since the results in
%implementation 1 are not very promising, we have decided to consider G-LRSP instead of
%SPP-LRSP.
In SPP-LRSP, each $z$ variable corresponds to a possible vehicle schedule, and since
all vehicles are identical, the vehicles are not differentiated with an index.
In the light of this information, to be able to relate a generated cut to model SPP-LRSP,
we project the variables in G-LRSP to a space with variables without the vehicle
index. Let $\sum_{h \in H_j}x_{kih} = u_{kij}$ for all $(k,i) \in A$, $j \in J$, and
$w_j = \sum_{h \in H_j}v_h$ for all $j \in J$.
Variable $u_{kij}$ will take value 1 if any vehicle located at facility $j$
travels link  $(k,i)$. $w_j$ is the number of vehicles used at facility $j$.
We start the implementation DISJCUTGEN by choosing the following graph based formulation
in step 1.
\begin{optprog}
{\bf [$M_2$]} & \sum_{j\in J}\sum_{k \in V} u_{kij} & = & 1 &  \forall  i \in I, \label{m2-c1}\\
   & -\sum_{i \in I} D_i \sum_{k\in V} u_{kij} + C^F_j t_j & \geq & 0 & \forall j \in J, \label{m2-c2}\\
   & \sum_{k\in V}u_{ikj}-\sum_{k \in V}u_{kij} & = & 0 & \forall i \in I, j \in J, \label{m2-c3}\\
   & TL w_j - \sum_{k\in V}\sum_{i \in V}t_{ki}u_{kij} & \geq & 0 & \forall j \in J, \label{m2-c4}\\
   & 0 \leq t_j  & \leq & 1 & \forall j \in J, \label{m2-c5} \\
   & 0 \leq u_{ikj} & \leq & 1 & \forall i,k \in V, j \in J, \label{m2-c6}\\
   & w_j & \geq & 0 & \forall j \in J. \label{m2-c7}
\end{optprog}
Constraints in formulation $M_2$ are similar to formulation G-LRSP. Constraints (\ref{m2-c1})
require that each customer must be served. Constraints (\ref{m2-c2}) are facility capacity
constraints. Constraints (\ref{m2-c3}) equate the number of vehicles in and out of a customer. 
G-LRSP has a set of constraints that are similar to (\ref{m2-c3}) but for each vehicle,  
$M_2$ has the aggregated version (aggregated for all vehicles located at the same facility). 
Again, constraints  (\ref{m2-c4}) are the aggregated version of the time limit constraints in G-LRSP. 
While constructing formulation $M_2$, in addition to projecting $(x,v)$ in G-LRSP
to $(u,w)$ (i.e.,aggregating vehicles located at the same facility),
we relaxed some constraints of G-LRSP such as the constraints that are including flow variables.

Flow variables ($y$ in G-LRSP) keep the amount of vehicle load
on the links. The reason for relaxing the constraints including flow variables is that we want to 
avoid generating cuts which are including flow variables (or any projection of flow variables) with 
nonzero coefficients. Based on the DWD procedure, we know how to rewrite the flow variables in terms 
of master problem variables ($z$) using relation (\ref{TotalYintermsTeta1}). 
%When we rewrite the generated cut in terms of the RMP variables $(t,z)$,
However, we do not know how to incorporate an information passed from the flow variables into the pricing 
solution algorithm. Therefore, we want to avoid generating a cut that includes flow variables.   
By using model $M_2$, we aim to generate disjunctive cuts that do not include flow variables.
%The reason for this relaxation is
%that we do not know how to incorporate a cut including flow variables ($f$ in G-LRSP).
%to eliminate the difficulty in solving CGLP.
Model $M_2$ enables a search for disjunctive cuts in the following form:
\begin{optprog}
& \sum_{j \in J}\alpha^t_{j}t_j + \sum_{j \in J}\alpha^w_{j}w_j + \sum_{j \in J}\sum_{k\in V}\sum_{i \in V}\alpha^u_{kij}u_{kij} & \geq & \beta. \label{cut3}
\end{optprog}
As in implementation 1, let $(\bar{t}, \bar{z}) \in {\mathbb{R}^{+}}^{|J|+\sum_{j \in J}|\bar{P}_{j}|}$
be a fractional solution to optimal LP relaxation of RMP.
We can find $(\bar{t}, \bar{w}, \bar{u}) \in {\mathbb{R}^{+}}^{|J|+|J|+|J|\times|A|}$
using relations $\bar{w}_j = \sum_{p \in \bar{P}_j}\bar{z}_p$,
$\bar{u}_{kij}  =  \sum_{p \in \bar{P}_j}l_{kip}\bar{z}_{p}$
for all $(k,i) \in A$, $j \in J$ where $l_{kip}$ is 1 if
pairing $p$ includes link $(k,i)$ and 0 otherwise.

If $(\bar{t}, \bar{w}, \bar{u})$ is fractional, then we consider the following disjunctions
in order to determine two disjunctive inequalities  violated by this solution.
\begin{enumerate}
\item A facility is either open or closed.
\begin{optprog}
& t_{j} \leq 0  & \vee & t_{j} \geq 1 & \forall j \in J. \label{imp2-disj1}
\end{optprog}
\item Total number of vehicles used at a facility is an integer.
\begin{optprog}
& w_{j} \leq  g_j & \vee & w_{j} \geq g_j +1 & \forall j \in J, \label{imp2-disj2}
\end{optprog}
for some integer $g_j$ and for all  $j \in J$.
\item Number of links into a set of customers is integer. Let $S$ be a customer
set such that $2 \leq |S| \leq 3$:
\begin{optprog}
& \sum_{j \in J}(\sum_{k \in V \backslash S}\sum_{i \in S} u_{kij}  & + &  \sum_{k \in V \backslash S}\sum_{i \in S}u_{ikj})& \leq & g \nonumber\\
& & \vee & \label{imp2-disj3} \\
& \sum_{j \in J}(\sum_{k \in V \backslash S}\sum_{i \in S} u_{kij}  & + & \sum_{k \in V \backslash S}\sum_{i \in S}u_{ikj}) & \geq & g + 1, \nonumber
\end{optprog}
%for some set of customers, $S$ such that $2 \leq |S| \leq 3$,
%$\sum_{j \in J}\sum_{k \in V \backslash S}\sum_{i \in S} \bar{u}_{kij} = b$
%\begin{optprog}
%& \sum_{j \in J}\sum_{k \in V \backslash S}\sum_{i \in S} \bar{u}_{kij} & = & b,
%\end{optprog}
%where $b$ is fractional and $g = 2\lfloor b \rfloor$.
for some integer $g$.
\end{enumerate}
%\zaupdate{mention how to choose disjunctions?}

By solving the CGLP constructed using model $M_2$ and one set of violated disjunctive
inequalities, we investigate for a cut that is maximally violated (by
fraction solution $(\bar{t}, \bar{w}, \bar{u})$) and is
in the form of (\ref{cut3}). A solution to the CGLP provides coefficients
$(\alpha, \beta) \in {\mathbb{R}}^{|J|+|J|\times |A|+1}$, and we can rewrite the cut
in terms of variables $t$, $w$ and $z$:
\begin{optprog}
& \sum_{j \in J}\alpha^t_{j}t_j + + \sum_{j \in J}\alpha^w_{j}w_j + \sum_{j \in J}\sum_{p \in P_j}\sum_{(i,k) \in A}\alpha^u_{ikj}l_{ikp}z_{p} \geq \beta.  \label{disj_cut_imp2}
\end{optprog}
Note that $w_j$ is the number of vehicles located at facility $j \in J$ and defined
with relation $\bar{w}_j = \sum_{p \in P_j}z_p$ in AUG-SPP-LRSP.
A cut in the form of (\ref{disj_cut_imp2}) can be added to the column generation.
Let $\omega$ be the dual variable corresponds to (\ref{disj_cut_imp2}).
In the pricing algorithm for facility $j$  for all $j \in J$, we update the cost of
link $(i,k)$ with $\alpha^u_{ikj}\omega$ for all $(i,k) \in A$ and we update the total cost 
of all labels at the end with fixed cost $\alpha^w_{j}\omega$.
%\zaupdate{Results of implementation 2? Effective? improment over implementation 1?}

\subsection{Implementation 3}
\label{sec:imp3-disjcut}
%\zaupdate{Since the performance of implementation 2 in which a relaxation of G-LRSP is used is not very good,}
In the third implementation, instead of a relaxation or a projection, model G-LRSP is
selected in step 1 of the disjunctive cut generation algorithm.
The CGLP will be formulated based on the LP relaxation of the G-LRSP. Recall that the model
is as follows:
\begin{optprog*}
{\bf [$M_3$]}  & \sum_{h\in H}\sum_{k \in N} x_{ikh} & = & 1 &  \forall  i \in I, \\
   & \sum_{k\in N} x_{ikh}-\sum_{k \in N} x_{kih} & = & 0 & \forall i \in N, h \in H, \\
   & \sum_{h \in H_j}\sum_{k \in I} y_{jkh} - L^F_j t_{j} & \leq & 0 &
       \forall j \in J, \\
   & y_{ikh} - L^V x_{ikh} & \leq & 0 & \forall (i,k) \in A, h \in H, \\
   & \sum_{k \in N} y_{ikh} - \sum_{k \in N} y_{kih} + D_i\sum_{k \in N}x_{ikh} & = & 0 & \forall i \in I, h \in H, \\
   & \sum_{(i,k) \in A} T_{ik} x_{ikh} - L^T v_h & \leq & 0 & \forall h \in H, \\
   & 0 \leq t_j  & \leq & 1 & \forall j \in J, \\
   & 0 \leq x_{ikh} & \leq & 1 & \forall i,k \in V, h \in H, \\
   & 0 \leq v_h   & \leq & 1 & \forall h \in H, \\
   & y_{ikh} & \geq & 0 &  \forall i,k \in V, h \in H.
\end{optprog*}
Based on model $M_3$ and one of the disjunctions (\ref{imp2-disj1}) - (\ref{imp2-disj3}), we
formulate the CGLP. Disjunction (\ref{imp2-disj3}) can be rewritten in terms of $x$ variables
in $M_3$:
\begin{optprog}
& \sum_{j \in J}(\sum_{k \in V \backslash S}\sum_{i \in S}\sum_{h \in H_j} x_{kih}  & + &  \sum_{k \in V \backslash S}\sum_{i \in S}\sum_{h \in H_j} x_{kih})& \leq & g \nonumber\\
& & \vee & \label{imp3-disj3} \\
& \sum_{j \in J}(\sum_{k \in V \backslash S}\sum_{i \in S}\sum_{h \in H_j} x_{kih}  & + & \sum_{k \in V \backslash S}\sum_{i \in S}\sum_{h \in H_j} x_{kih}) & \geq & g + 1, \nonumber
\end{optprog}
for some set of customers, $S$.
%such that $2 \leq |S| \leq 3$,
%$\sum_{j \in J}\sum_{k \in V \backslash S}\sum_{i \in S} \sum_{h \in H_j} \bar{x}_{kih} = b$
%where $b$ is fractional and $g = 2\lfloor b \rfloor$.
The CGLP based on model $M_3$ enables a search for violated disjunctive cuts in the following form:
\begin{optprog}
    & \sum_{j \in J} \alpha^t_{j} t_j + \sum_{h \in H}\alpha^v_{h} v_h + \sum_{(k,i)\in A}\sum_{h \in H}\alpha^y_{kih}y_{kih} + \sum_{(k,i)\in A}\sum_{h \in H}\alpha^x_{kih}x_{kih} & \geq & \beta. \label{imp3-c1}
\end{optprog}
Since we do not favor cuts with nonzero coefficients for flow variable vector $y$, we add
the following constraint to the CGLP to force the coefficients to be zero:
\begin{optprog}
&\alpha^y_{kih} & = & 0 & \forall  (i,k) \in A, h \in H. \label{0coef-flow}
\end{optprog}
%In CGLP, we force constraint (\ref{0coef-flow}) to set the coefficients for flow variables
%to zero because of the reason explained in implementation 2.
We let $(\bar{t}, \bar{z}) \in {\mathbb{R}^{+}}^{|J|+\sum_{j \in J}|\bar{P}_{j}|}$
be a fractional solution to the LP relaxation of RMP where $\bar{P}_{j}$ be the 
set of pairings with nonzero values in the solution, i.e.,$\bar{z}_p > 0$ for all
$p \in \bar{P}_{j}$ for all $j \in J$. We can generate vector
$(\bar{t}, \bar{v}, \bar{x}, \bar{y}) \in {\mathbb{R}^{+}}^{|J|+|J|\times|H|+|H|\times|A| + |H|\times|A|}$
representing the projection of the RMP solution for model $M_3$. Since we have identical vehicles,
we can assign nonzero pairings from 1 to $|\bar{P}_j|$ for all $j \in J$ to vehicles (in $M_3$) arbitrarily, 
starting from the first vehicle. Let $(h \leftrightarrow p)$ represent such correspondence where
$p$ is the $h^{st}$ smallest index in $\bar{P}_j$ for all $h = 1, ..,|\bar{P}_j|$
and $j \in J$. Note that a vector  $(\bar{t}, \bar{z})$ may lead to multiple
corresponding $(\bar{t}, \bar{v}, \bar{x}, \bar{y})$ vectors under the identical
vehicles assumption. By defining correspondence $(h \leftrightarrow p)$ between
$h$ and $p$, we restrict the correspondence between vectors  $(\bar{t}, \bar{z})$
and $(\bar{t}, \bar{v}, \bar{x}, \bar{y})$ to one-to-one.
To construct $(\bar{t}, \bar{v}, \bar{x}, \bar{y})$, we use the following relations
based on defined correspondence between index $h$ and $p$:
\begin{eqnarray*}
 \bar{v}_h & = & \left \{ \begin{array}{ll}
                 \bar{z}_p & \mbox{If $(h \leftrightarrow p)$, $\forall p \in \bar{P}_j, h \in H_j, j \in J$,} \\
                 0         & \mbox{If $h = |\bar{P}_j|+1, .. , |H_j|$, $\forall j \in J$,}
                 \end{array}  \right.      \\
 \bar{x}_{ikh} & = & \left \{ \begin{array}{ll}
                 l_{ikp}\bar{z}_p & \mbox{If $(h \leftrightarrow p)$, $\forall p \in \bar{P}_j, h \in H_j, j \in J$,} \\
                 0         & \mbox{If $h = |\bar{P}_j|+1, .. , |H_j|$, $\forall j \in J$,}
                 \end{array}  \right. \\
 \bar{y}_{ikh} & = & \left \{ \begin{array}{ll}
                 f_{ikp}\bar{z}_p & \mbox{If $(h \leftrightarrow p)$, $\forall p \in \bar{P}_j, h \in H_j, j \in J$,} \\
                 0         & \mbox{If $h = |\bar{P}_j|+1, .. , |H_j|$, $\forall j \in J$,}
                 \end{array}  \right.
\end{eqnarray*} %{optprog}
where $l_{ikp}$ is 1 if pairing $p$ includes link $(i,k)$ and $f_{ikp}$
is equal to the vehicle load on link  $(i,k)$ of pairing $p$, for all
$(i,k) \in A$, $p \in P_j$ and  $j \in J$.
After a cut that is in the form of (\ref{imp3-c1}) and satisfying (\ref{0coef-flow}) is 
generated for a given fractional solution, $(\bar{t}, \bar{v}, \bar{x}, \bar{y})$,  
we can write more cuts by circulating coefficients $\alpha^x_{ikh}$ and
$\alpha^v_h$ over all $h \in H_j$ for all $j \in J$ using the assumption of identical
vehicles. In other words, we assume that we obtain multiple $(\bar{t}, \bar{v}, \bar{x}, \bar{y})$
vectors from one  $(\bar{t}, \bar{z})$ vector (one-to-many correspondence between solution vectors).
%Note that there are many other possible
%mappings between vehicle indices and pairings other than assigning pairings to vehicles starting
%with the smallest index. %We could find some of these mappings by circulatinf
Then, we can combine the coefficients,
$\alpha^x$ and $\alpha^v$ indexed by $h$ (for same facility), and rewrite the cut (\ref{imp3-c1}):
\begin{optprog}
 & \objective{\sum_{j \in J} \alpha^t_{j} t_j + \sum_{j \in J}\alpha^{v_c}_{j}\sum_{h \in H_j}v_h + \sum_{j \in J}\sum_{(k,i)\in A}\sum_{h \in H_j}\alpha^{x_c}_{kij}x_{kih} \geq \beta^c} \label{imp3-c2} \\
%such that  &  \nonumber \\  %\mbox{such that} \nonumber \\
such that  & \alpha^{v_c}_{j} & = & \sum_{h \in H_j}\alpha^v_{h}, \label{imp3-con1}\\
 & \alpha^{x_c}_{ikj} & = & \sum_{h \in H_j}\alpha^x_{ikh}, \label{imp3-con2} \\
 & \beta^c & = & |I|\beta, \nonumber
\end{optprog}
where $|H_j| = |I|$ for all $j \in J$. Then, we can project the cut to the RMP space:
\begin{optprog}
& \sum_{j \in J}\alpha^t_{j}t_j + \sum_{j \in J}\alpha^{v_c}_{j}w_j +
\sum_{j \in J}\sum_{(i,k) \in A}\alpha^{x_c}_{ikj}\sum_{p \in P_j}l_{ikp}z_{p} \geq \beta^c  \label{disj_cut_imp3}
\end{optprog}
As a summary, We solve the CGLP to find a cut in the form of (\ref{imp3-c1}) that violates the current
fractional solution. Then we update the coefficients of (\ref{imp3-c1}) to get
cut in the form of (\ref{imp3-c2}). Since the CGLP is solved for  (\ref{imp3-c1})
and not for (\ref{imp3-c2}), the fractional solution $(\bar{t}, \bar{z})$ may not violate 
(\ref{disj_cut_imp3}).
%However, a fractional solution violating cut (\ref{imp3-c1}) may not violate (\ref{imp3-c2}) since we
%are updating coefficients after we solve the CGLP.
Therefore, we designed implementation 4 to be able to overcome this problem.

\begin{comment}
\subsubsection{Method 1.}
\begin{enumerate}
\item A disjunctive cut in the following form is genetated:
  \begin{optprog}
    & \sum_{j \in J} \alpha_{t_j} t_j + \sum_{h \in H}\alpha_{y_h} y_h + \sum_{k\in V}\sum_{i \in V}\sum_{h \in H}\alpha_{kih}x_{kih} & \geq & \beta \label{cut1}
  \end{optprog}
Note that the coefficients for flow variables, $f$ are set to zero in (\ref{cut1}).
\item  The cut generating LP is constructed using the formulation EB-LRSP. In the cut generating LP, the $\alpha$
coefficients for variable  $f$ are fixed  to be zero.
\item After a cut in form (\ref{cut1}) is generated, using the assumption of identical vehicles and
symmetric solutions, coefficients indexed with $h$ are combined to get the following cut:
  \begin{optprog}
    & \sum_{j \in J} \alpha_{t_j} t_j + \sum_{j \in J}\alpha^c_{y_j}\sum_{h \in H_j} y_h + \sum_{k\in V}\sum_{i \in V}\sum_{j \in J}\alpha^{c_x}_{kij}\sum_{h \in H_j}x_{kih} & \geq & \beta^c \label{cut2}
  \end{optprog}
\end{enumerate}
\end{comment}

\subsection{Implementation 4}
This implementation is designed to improve implementation 3. Similarly, step 1 of DISJCUTGEN
chooses model $M_3$ and the CGLP is defined based on this model.
%CGLP based on model $M_3$.
However, we directly search for the most violated cut in the form of (\ref{imp3-c2})
instead of the form (\ref{imp3-c1}). To construct a cut of the form (\ref{imp3-c2}),
we force constraints (\ref{0coef-flow}) in the CGLP, i.e., $\alpha$ coefficients for 
flow variables (for $y$) are set to zero. In addition, to be able to
obtain the combined (aggregated for all vehicles located at the same facility) coefficients
for $x$ and $v$, we set $\alpha$ coefficients for variables $x_{ikh}$ and $v_h$ to be equal for all $h \in H_j$
for $j \in J$. 
\begin{optprog}
 & \alpha^{v_c}_{j} & = & \alpha^v_{h}, & \forall h \in H_j, \forall j \in J,\label{imp4-con1}\\
 & \alpha^{x_c}_{ikj} & = & \alpha^x_{ikh} & \forall h \in H_j, \forall j \in J. \label{imp4-con2} 
\end{optprog}
As in  implementation 3, we use one of the disjunctions (\ref{imp2-disj1}),
(\ref{imp2-disj2}) and (\ref{imp3-disj3}) to obtain a pair of disjunctive 
inequalities.

%\zaupdate{Results of implementation 4?}


\section{Cutting Planes Based on Arc Based Formulations}
Another approach to derive problem specific cutting planes for LRSP
is to study the LRSP polytope or its relaxations based on arc based
formulations (like G-LRSP or similar relaxed models) that are defined
based on variables that correspond to links between the nodes in the problem.
When a cut is derived in terms of link variables, then this cut can be
projected to SPP-LRSP by defining certain relationships between variables.

The literature does not include any studies devoted to polyhedral
structure of LRSP. However, there are studies that investigate
cuts and valid inequalities for VRP and LRP which are special cases
of LRSP. In this section, we investigate some sets of cutting
planes that are derived for VRP or LRP and originally defined in terms of
variables corresponding to edges in the problem. These cutting planes are
also valid for LRSP. Here, we discuss how we can incorporate these inequalities
to SPP-LRSP and whether these inequalities have a positive impact to the model.
In fact, we learned from the literature or showed that most of these inequalities 
are already implied by SPP-LRSP.

It is common to use two index vehicle flow formulations in the studies
for VRP and LRP that investigate cutting planes. These formulations use binary
variables, defined for each link, representing whether the link is used in the
solution or not. Cutting planes based on link variables can easily be
modified for SPP-LRSP and may be used to extend the branch and price
algorithm. Let $x_{ik}$ be a link variable used in VRP or LRP
formulations where $x_{ik}$ equals to 1 if link $(i,k)$ is used in
the solution. When we have a cut in terms of $x$, we can use the  relation,
\begin{optprog}
& x_{ik} & = & \sum_{j \in J}\sum_{p \in P_j}l_{ikp}z_p & \forall (i,k) \in A, \label{xz-convert}
\end{optprog}
%$x_{ik}  = \sum_{j \in J}\sum_{p \in P_j}l_{ikp}z_p$
to project the cut into the RMP space where $l_{ikp} = 1$ if link $(i,k)$ is in
pairing $p$. In some graph based formulations, like G-LRSP, variables $x_{ik}$ are
replaced with three index variables that also include a vehicle index, such as
$x_{ikh}$, where $x_{ikh}$ is 1 if vehicle $h \in H$ travels link $(i,k)$ in a solution
to a three index formulation. We know the relation between a two index and a three
index variable, which is $x_{ik} = \sum_{h \in H}x_{ikh}$.
%In addition, some formulations the cutting planes based
%- Cutting planes based on link variables can easily be adapted to the pricing problem.
%- Sum over all vehicles to obtain two index variables
%- some inequalities are implied by SP-LRSP based on ESPPRC pricing problem.
%\begin{optprog}
%& x_{ik} & = & \sum_{j \in J}\sum_{p \in P_j}b_{ikp}z_p & \forall i \in N, k \in N \label{xz-convert}
%\end{optprog}

Some cutting planes used in branch and cut algorithms for a VRP or LRP
are implied by the set partitioning formulations of these problems.
In the literature, there are limited number of studies investigating
the theoretical relationship between a set partitioning formulation and a
graph based formulations of the VRP, but there are none for the LRP.
\citet{LetchfordG06} theoretically prove that a solution to the LP relaxation
of a set partitioning formulation for the VRP (with elementary resource constrained
paths) implies some of the valid inequalities (such as {\em knapsack large multistar,
sub tour elimination, generalized large multistar, hypotour-like inequalities})
that are derived for a graph based formulation in the two index space.
This result is also applicable to the LRSP as the LP relaxation of SPP-LRSP
satisfies these type of cuts. \citet{Ropke07} do a similar study for the VRPTW with
pickup and delivery; they investigate the relationship between a branch and cut algorithm
and a branch, cut and price algorithm for the problem. They show that some of the
valid inequalities for the problem (e.g. {\em fork inequalities, reachability inequalities,
precedence inequalities}) are implied by the set partitioning formulation of the problem
with a pricing problem that generates elementary resource constrained paths.

We started our study by investigating two of the most common classes of inequalities
for the VRP:
%we considervehicle capacity cuts,
fractional capacity inequalities and rounded capacity inequalities.
%which we are also valid for LRSP.

\paragraph{Fractional Capacity Inequalities}
\citet{LetchfordG06} note that fractional capacity inequalities are
dominated by generalized large multistar (GLM) inequalities and GML inequalities are
included in knapsack large multistar inequalities, which are proved to be implied by
the set partitioning formulation of the VRP. Therefore, \citet{LetchfordG06}'s
results prove that fractional capacity inequalities are redundant for the set
partitioning formulation. Here, we re-derive this result.
Let $S$ be a set of customers, then the fractional capacity
inequality for $S$ based on two index variables is as follows:
\begin{optprog}
& \sum_{i \in S}\sum_{k \in V\backslash S}x_{ik} & \geq & \frac{\sum_{i \in S}D_i}{L^V}. \label{frac-cap-inq}
\end{optprog}
In other words, total capacity of the vehicles leaving $S$ must be greater than the
total demand of customers in $S$. Using relation (\ref{xz-convert}), we can
rewrite inequality (\ref{frac-cap-inq}) in terms of $z$ variables.
\begin{optprog}
& \sum_{j \in J}\sum_{p \in P_j}\sum_{i \in S}\sum_{k \in V\backslash S}l_{ikp}z_p
= \sum_{j \in J}\sum_{p \in P_j}\alpha_{p}z_p & \geq & \frac{\sum_{i \in S}D_i}{L^V}, \label{frac-cap-inq2}
\end{optprog}
where $\alpha_{p}$ is the coefficient that corresponds to $z_p$ for all
$p \in \cup_{j \in J}P_{j}$. $\alpha_p$ is the number of links that are
from set $S$ to set $V\backslash S$ and included in pairing $p$.
%Let $p$ be a pairing such that $p \in \cup_{j \in J}P_{j}$ where $P_j$
%is the set of elementary feasible pairings for facility $j$, for all $j \in J$.
Let $I_p$ be the set of customers visited by pairing $p$, i.e.,
$I_p=\{i \in I | a_{ip} = 1 \}$. For ease of notation, we define $n_K$ as the
minimum number of vehicle routes that are needed to cover customer demand in set
$K$, in other words, $n_K = \lceil\frac{\sum_{i \in K}D_i}{L^V}\rceil$. Since a
pairing represents a feasible vehicle schedule, by construction, the total number of 
links (in some $p$) leaving customers in set $K$ that are visited by $p$ 
($K \subseteq I_p$) is at least $n_K$. The following lemma is stated based on this fact.
%Then, the following is true from the feasibility of pairings by construction.
\begin{lemma}
\label{contr-p-fc}
Let $K$ be the intersection of set of customers visited by $p$ with $S$ in
inequality (\ref{frac-cap-inq}), i.e.,$K = I_p \cap S$. If $K \neq \emptyset$,
then the contribution (coefficient) of variable $z_p$ to the left hand side of
the inequality (\ref{frac-cap-inq2}) is at least $n_K$. In other words, 
$\alpha_p \geq n_K$.
%($\alpha_p \geq n_K$).
\end{lemma}
Let $(t^*,z^*) \in {\mathbb{R}^{+}}^{|J|+|\cup_{j \in J}P^*_{j}|}$
be a solution to the LP relaxation of SPP-LRSP where $P^*_{j}$ is the current set
of pairings for facility $j \in J$.
\begin{prop}
If a vector $(t^*,z^*)$ satisfies (\ref{spp_con1}) and (\ref{spp_con2})
then it also satisfies (\ref{frac-cap-inq2}).
%vector $x^*$ transformed using (\ref{xz-convert})
%satisfies the fractional capacity inequalities (\ref{frac-cap-inq}).
\end{prop}
\begin{proof}
%Let $I_p=\{i \in I | a_{ip} = 1 \}$ be the set of customers visited
%by $p$ and $K_p = I_p \cap S$ for all $p \in \cup_{j \in J}P^*_{j}$.
We multiply constraint (\ref{spp_con1}) for each customer $i$ with its
demand $D_i$ for all $i \in S$. Then, we take the summation of the obtained
constraints for all customers in set $S$:
\begin{optprog}
& \sum_{j \in J}\sum_{p \in P^*_j}\sum_{i \in S}D_ia_{ip}z^*_{p} & = & \sum_{i \in S}D_i, \\
& \sum_{j \in J}\sum_{K \subseteq S}\sum_{p \in P^*_{Kj}}D(K) z^*_{p} & = & D(S),  \label{sumsetdemS}
\end{optprog}
where $D(K) = \sum_{i \in K}D_i$ and $P^*_{Kj}$ is the set of pairings visiting
all customers in set $K$ and possibly some other customers in $I\backslash S$
for all nonempty subsets of $S$, i.e. $I_p \cap S = K$ for all $p \in P^*_{Kj}$.
$P^*_{Kj}$ is defined for all $K \subseteq S$ and for all $j \in J$.
Note that $P^*_{K_1j} \cap P^*_{K_2j} = \emptyset$ for all $K_1, K_2 \subseteq S$,
$K_1 \neq K_2$ and $K_1$, $K_2$ are nonempty. Since the pairings are feasible with respect to vehicle capacity
from construction, $n_{K}L^V \geq D(K)$ for all $K \subseteq S$. Therefore,
\begin{optprog}
%& \sum_{j \in J}\sum_{p \in P_j}\sum_{i \in S}D_ia_{ip}z^*_{p} & = & \sum_{i \in S}D_i, \\
%& \sum_{j \in J}\sum_{K \subseteq S}\sum_{p \in P^*_{Kj}}D(K) z^*_{p} & = & D(S),  \label{sumsetdemS} \\
& \sum_{j \in J}\sum_{K \subseteq S}\sum_{p \in P^*_{Kj}}n_{K}L^V z^*_{p} \geq  \sum_{j \in J}\sum_{K \subseteq S}\sum_{p \in P^*_{Kj}}D(K) z^*_{p} & = & D(S),  \label{sumsetdemS2} \\
& \sum_{j \in J}\sum_{K \subseteq S}\sum_{p \in P^*_{Kj}}n_{K} z^*_{p} & \geq & \frac{D(S)}{L^V}
\end{optprog}
From lemma \ref{contr-p-fc}:
\begin{optprog}
%& \sum_{i \in S}\sum_{k \in V\backslash S} x^*_{ik} & \geq & \sum_{j \in J}\sum_{K \subseteq S}\sum_{p \in P^*_{K}}n_K z^*_{p} \geq \frac{D(S)}{L^V}
& \sum_{j \in J}\sum_{K \subseteq S}\sum_{p \in P^*_{Kj}}\alpha_p z^*_p & \geq & \sum_{j \in J}\sum_{K \subseteq S}\sum_{p \in P^*_{K}}n_K z^*_{p} & \geq & \frac{D(S)}{L^V}.
\end{optprog}
Therefore, inequality (\ref{frac-cap-inq2}) is satisfied for any solution of the LP relaxation
of the SPP-LRSP.
\qed
\end{proof}

\paragraph{Rounded Capacity Inequalities}
Another common class of inequalities is rounded capacity inequalities. These
inequalities dominate the fractional capacity inequalities and they are
in the following form:
\begin{optprog}
& \sum_{i \in S}\sum_{k \in V\backslash S}x_{ik} & \geq & n_S,\label{rnd-cap-inq}
\end{optprog}
where $n_S = \lceil\frac{\sum_{i \in S}D_i}{L^V}\rceil$ for some customer set $S$.
Neither \citet{LetchfordG06} nor \citet{Ropke07} has a result that says
that the rounded capacity inequalities are implied by set partitioning formulation.
Here, we investigate whether these inequalities can make an impact to the SPP-LRSP.
Based on relation (\ref{xz-convert}), if we rewrite (\ref{rnd-cap-inq}) in terms of
$z$ variables, we obtain an inequality in the form of (\ref{frac-cap-inq2}) with
a right hand side that is equal to $n_S$.

Let $(t^*,z^*) \in {\mathbb{R}^{+}}^{|J|+|\cup_{j \in J}P^*_{j}|}$
be a solution to the LP relaxation of SPP-LRSP where $P^*_{j}$ is the current set
of pairings for facility $j \in J$.
\begin{lemma}
If $n_S = 1$ or $n_S = |S|$, then $(t^*,z^*)$ satisfies inequality
(\ref{rnd-cap-inq}).
\end{lemma}
\begin{proof}
If $n_S = 1$,  inequality (\ref{rnd-cap-inq}) is equal to a
subtour elimination inequality, which is implied by the set partitioning formulation
(proved in \citet{LetchfordG06}).
% or from the constraint (\ref{spp_con1}) for any customer in set $S$.

If $n_S = |S|$, then no two customers in $S$ can be visited in the same route.
Thus, if Lemma \ref{contr-p-fc} is applied, the coefficient of
pairing $p$ that visits customers in set $K \subseteq S$ in (\ref{rnd-cap-inq})
is at least $n_K = |K|$. If we take the summation of constraints (\ref{spp_con1})
for all customers in set $S$, the validity of inequality (\ref{rnd-cap-inq}) follows:
\begin{optprog}
& \sum_{j \in J}\sum_{p \in P^*_j}\sum_{i \in S}a_{ip}z^*_{p} & = & |S|, \\
& \sum_{j \in J}\sum_{K \subseteq S}\sum_{p \in P^*_{Kj}}\alpha_p z^*_p  \geq \sum_{j \in J}\sum_{K \subseteq S}\sum_{p \in P^*_{Kj}}|K|z^*_{p} & = & |S|.
%& \sum_{i \in S}\sum_{k \in V\backslash S} x^*_{ik} \geq \sum_{j \in J}\sum_{K \subseteq S}\sum_{p \in P^*_{Kj}}|K|z^*_{p} & = & |S|.
\end{optprog}
\qed
\end{proof}
If the minimum number of vehicle routes that are needed to serve customers in $S$
is greater than 1 and less than $|S|$, then
%$1 < n_S < |S|$, then
the rounded capacity inequalities may be violated by some fractional LP solutions of
SPP-LRSP. An example of such a solution for
a small instance follows:
\begin{example}
{\bf An instance with 4 customers and 2 facilities:}
\begin{itemize}
\item Let $I = \{1,2,3,4\}$, $J = \{\mbox{facility}_1,\mbox{facility}_2\}$.
\item Demand is 20 for all customers. Vehicle capacity is 50, and facility capacity is 70.
\item Let $(t^*,z^*)$ be the LP relaxation solution value at root node for SPP-LRSP.
\item Nonzero variable values and the pairings represented by pairing variables are:
%\begin{eqnarray*}
\begin{optprog*}
& t^*_{\mbox{facility}_1} = t^*_{\mbox{facility}_2}  & =  & 1 \\
& z^*_1 & = & 1   & : & \mbox{facility$_1$} - 4 - \mbox{facility$_1$} \\
& z^*_2 & = & 0.5 & : &  \mbox{facility$_2$} - 3 - 1 - \mbox{facility$_2$} \\
& z^*_3 & = & 0.5 & : &  \mbox{facility$_2$} - 1 - 2 - \mbox{facility$_2$} \\
& z^*_4 & = & 0.5 & : &  \mbox{facility$_2$} - 3 - 2 - \mbox{facility$_2$}
\end{optprog*}
\item Rounded capacity inequality for $S = \{1,2,3\}$ is violated,
\begin{optprog*}
& \sum_{i \in S}\sum_{k \in V\backslash S}x^*_{ik} = \sum_{j \in J}\sum_{p \in P_j}\sum_{i \in S}\sum_{k \in V\backslash S}l_{ikp}z^*_p & = & 1.5 < 2
\end{optprog*}
\end{itemize}
\end{example}
At this point of the research, we have decided to leave the implementation
of these inequalities in our branch and price algorithm as a future work.
%\zaupdate{Can I code separation for rounded capacity inequalities?}

Next, we investigate some cutting planes used for LRP, and discuss their
validity for LRSP and their effect in formulation SPP-LRSP.

\paragraph{Chain Barring Inequalities}
\citet{Laporte86} introduced chain barring constraints to eliminate paths
from one depot to another in a solution of the LRP. Let $r=(j_1, i_1, i_2, i_3, .., i_k, j_2)$
be a path from node $j_1$ to $j_2$ where $j_1$ and $j_2$
are two depots and $i_1$, $i_2$, .., $i_k$ are customer nodes. To eliminate
path $r$, the following chain barring inequalities are defined:
\begin{optprog}
& x_{j_1i_1} + 3x_{i_1i_2} + x_{i_2j_2} & \leq &  4  & \mbox{if $k = 2$} \label{chb-1}\\
& x_{j_1i_1} + 2\sum_{t=1}^{k-1}x_{i_ti_{t+1}} + x_{i_kj_2} & \leq & 2k -1  & \mbox{if $k \geq 3$} \label{chb-2}
\end{optprog}
We will prove that the LP relaxation of SPP-LRSP satisfies the following chain barring
inequalities that are obtained by rewriting (\ref{chb-1}) and (\ref{chb-2}) in
terms of $z$ variables.
%Using relation (\ref{xz-convert}), inequalities
%(\ref{chb-1}) and (\ref{chb-2}) can be rewritten in SPP-LRSP variable
%space.
\begin{optprog}
& \sum_{p \in P_{j_1}}l_{j_1i_1p}z_p + \sum_{j \in J}\sum_{p \in P_j}l_{i_1i_2p}z_p + \sum_{p \in P_{j_2}}l_{i_2j_2p}z_p & \leq &  4  & \mbox{if $k = 2$} \label{chb-1b}\\
& \sum_{p \in P_{j_1}}l_{j_1i_1p}z_p + 2\sum_{j \in J}\sum_{p \in P_j}\sum_{t=1}^{k-1}l_{i_ti_{t+1}p}z_p + \sum_{p \in P_{j_2}}l_{i_kj_2p}z_p & \leq & 2k -1  & \mbox{if $k \geq 3$} \label{chb-2b}
%x_{j_1i_1} + 2\sum_{t=1}^{k-1}x_{i_ti_{t+1}} + x_{i_kj_2}
\end{optprog}
For the ease of notation, let $\alpha_p$ be the coefficient corresponds to
variable $z_p$ for all $p \in \cup_{j \in J}P_{j}$. Let $(t^*,z^*)$ be a solution
to the LP relaxation of SPP-LRSP and in this solution
$r=(j_1, i_1, i_2, i_3, .., i_k, j_2)$ be a path from
depot $j_1$ to $j_2$ including $k$ customers.
%We can define vector $x^*$
%using $z^*$ and relation (\ref{xz-convert}).
\begin{lemma}
\label{contr-p}
%Let $p \in \cup_{j \in J}P_{j}$ where $P_j$ is the set of elementary feasible
%pairings for facility $j$, for all $j \in J$.
If pairing $p \in \cup_{j \in J}P_{j}$ visits exactly $n$ customers
in path $r$, then it can visit at most $n$ arcs in path $r$.
%either the arc $(j_1,i_1)$ or $(i_k,j_2)$.
Depending on the number of customers in the path, the contribution (coefficient)
of $z^*_p$ in the left hand side of inequalities (\ref{chb-1b}) and (\ref{chb-2b})
is as follows:
\begin{itemize}
\item [i.] If $k = 2$, the contribution of $p$ in (\ref{chb-1b}) is $3(n-1)+1$.
\item [ii.] If $k \geq 3$, the contribution of $p$ in (\ref{chb-2b}) is $2(n-1)+1 = 2n-1$.
\end{itemize}
%it can visit at most 1 arc in path $r$, either the arc $(j_1,i_1)$
%or $(i_k,j_2)$.
\end{lemma}
\begin{proof}
If $p$ visits 1 customer in the path, then it can use at most one
arc in $r$. If $p$ only visits one of the customers in $\{i_2, ..,i_{k-1}\}$,
it cannot include any arc included in $r$. If $p$ visits
either $i_1$ or $i_k$, then it can include either the arc $(j_1,i_1)$ or
$(i_k,j_2)$. The path cannot include both arcs since from the construction
of $p$, $p$ can only visit one of the facilities. Therefore, the contribution of
variable $z^*_p$ to both (\ref{chb-1b}) and (\ref{chb-2b}) can be at most
1 because of either $(j_1,i_1)$ or $(i_k,j_2)$.
%$x_{j_1i_1}$ or $x_{i_kj_2}$.

If $p$ visits $n$ ($n>1$) customers included in $r$, then it can use the
largest number of arcs in $r$ only if it visits these customers consecutively
by using $n - 1$ arcs. And if $p$ visits the first ($i_1$) or last customer ($i_k$)
in the path, then it can also include either the arc $(j_1,i_1)$ or $(i_k,j_2)$.
%$x_{j_1i_1}$ or $x_{i_kj_2}$.
Therefore, $p$ can use at most $n$ arcs in $r$. The contribution of
$p$ to  (\ref{chb-1b}) and (\ref{chb-2b}) follows from the coefficients
in the inequalities and the fact that at most $n-1$ of $n$ arcs are between two customers.\qed
\end{proof}
\begin{prop}
If a vector $(t^*,z^*)$ satisfies (\ref{spp_con1}) and (\ref{spp_con2})
then it also
%vector $x^*$ transformed using (\ref{xz-convert})
satisfies the chain barring inequalities (\ref{chb-1b}) and (\ref{chb-2b}).
\end{prop}
\begin{proof}
Let $r=(j_1, i_1, i_2, i_3, .., i_k, j_2)$ be a path from
depot $j_1$ to $j_2$ including customers $i_1, i_2, .., i_k$
and assume that $z^*$ violates either (\ref{chb-1b}) or (\ref{chb-2b}).
Let $P^*_{nj}$ be the set of pairings visiting exactly $n$ of the customers
in $r$ for facility $j$, for all $j \in J$ and $n = 1, ..,k$.

We take the summation of constraints (\ref{spp_con1}) for customers in set
$S=\{i_1, i_2,.., i_k \}$:
\begin{optprog}
& \sum_{j \in J}\sum_{p \in P^*_j}\sum_{i \in S}a_{ip}z^*_{p} & = & k, \\
& \sum_{j \in J}\sum_{n = 1}^{k}\sum_{p \in P^*_{nj}}n z^*_{p} & = & k.  \label{sumset-k}
\end{optprog}
For $k = 2$, using lemma \ref{contr-p} part (i):
\begin{optprog}
& \sum_{j \in J}\sum_{p \in P^*_j}\alpha_p z_p & \leq & \sum_{j \in J}\sum_{n = 1}^{2}(3n-2)\sum_{p \in P^*_{nj}} z^*_{p} = \sum_{j \in J}\sum_{p \in P^*_{1j}}z^*_{p} + 4\sum_{j \in J}\sum_{p \in P^*_{2j}}z^*_{p} \nonumber\\
& & \leq & 2(\sum_{j \in J}\sum_{p \in P^*_{1j}}z^*_{p} + 2\sum_{j \in J}\sum_{p \in P^*_{2j}}z^*_{p}) = 4 \label{chain-2}
\end{optprog}
%\begin{optprog}
%& x^*_{j_1i_1} + 3x^*_{i_1i_2} + x^*_{i_2j_2} & \leq & \sum_{j \in J}\sum_{n = 1}^{2}(3n-2)\sum_{p \in P^*_{nj}} z^*_{p} = \sum_{j \in J}\sum_{p \in P^*_{1j}}z^*_{p} + 4\sum_{j \in J}\sum_{p \in P^*_{2j}}z^*_{p} \nonumber\\
%& & \leq & 2(\sum_{j \in J}\sum_{p \in P^*_{1j}}z^*_{p} + 2\sum_{j \in J}\sum_{p \in P^*_{2j}}z^*_{p}) = 4 \label{chain-2}
%\end{optprog}
The equality in (\ref{chain-2}) follows from the equality (\ref{sumset-k}) for $k=2$. Therefore,
(\ref{chb-1b}) cannot be violated for any fractional solution.

\noindent For $k \geq 3$, using lemma \ref{contr-p} part (ii):
\begin{optprog}
&\sum_{j \in J}\sum_{p \in P^*_j}\alpha_p z_p & \leq &  \sum_{n = 1}^{k}(2n-1)\sum_{j \in J}\sum_{p \in P^*_{nj}}z^*_{p}  \nonumber\\
& & \leq & \sum_{n = 1}^{k}(n \frac{(2k-1)}{k})\sum_{j \in J}\sum_{p \in P^*_{nj}}z^*_{p} & = & 2k - 1 \label{chain-k}
\end{optprog}
%\begin{optprog}
%& x^*_{j_1i_1} + 2\sum_{t=1}^{k-1}x^*_{i_ti_{t+1}} + x^*_{i_kj_2} & \leq &  \sum_{n = 1}^{k}(2n-1)\sum_{j \in J}\sum_{p \in P^*_{nj}}z^*_{p}  \nonumber\\
%& & \leq & \sum_{n = 1}^{k}(n \frac{(2k-1)}{k})\sum_{j \in J}\sum_{p \in P^*_{nj}}z^*_{p} & = & 2k - 1 \label{chain-k}
%\end{optprog}
To second inequality in (\ref{chain-k}) follows from the fact that $(2n-1) \leq n\frac{(2k-1)}{k}$ 
for all $n = 1..k$. To obtain the equality,  multiply (\ref{sumset-k}) with $(2k-1)/k$. 
%he inequality  And, the last equality, 
The right hand side of (\ref{sumset-k}) after the multiplication with  $(2k-1)/k$ is equal to
$2k-1$. Therefore, (\ref{chb-2b}) can not be violated  for any fractional solution of SPP-LRSP. \qed
\end{proof}


\paragraph{Facility Capacity Inequalities}
\citet{Belenguer06} define the following inequalities to satisfy facility capacity constraints:
\begin{optprog}
& \sum_{i \in S \cup \{j\}}\sum_{k \in V\backslash (S\cup \{j\})} x_{ik}  & \geq & 1 &
       \forall j \in J, \forall S \subseteq I | \sum_{i \in I}D_i) > L^F_j  \label{FacCap2b}
\end{optprog}
\zaupdate{I do not think this is implied by AUG-SPP-LRSP. But, also I do not think, they
would be helpful. Since there are multiple customers visited by more than one facility in a 
fractional solution. This forces total flow (in the left hand side of inequality) to be $\geq$ 1.
So either, I should comment on that, or remove this inequality from the section.}

%The LP relaxation of the SPP-LRS does not imply these constraints since facility
%capacity constraints are not part of the pricing problem and integer feasibility of
%the constraints are not guaranteed either with pricing or master problem constraints.

%
%\zaupdate{lower bound for number of links from depots to customers based
%on vehicle capacity (rounded vehicle capacity ineqaulity for all customers}

%\zaupdate{Look at Jan1207/ convert a cut in x space to z space}

%\zaupdate{lower bound for $\sum_{j \in J}v_j >= ? $, see may1608}
\vspace{1in}
\zaupdate{Is there any other possible contribution to this chapter? Do you have any suggestions?}
